\documentclass[11pt]{ctexart}
\usepackage[a4paper,margin=1in]{geometry}
\usepackage{graphicx}
\usepackage{xcolor}
\usepackage{hyperref}
\definecolor{refgray}{gray}{0.45}
\title{\textbf{市场最新观点汇总}\\[0.4em]{\large 全球资产定价逻辑正在从“总量叙事”转向“结构性分化”，市场对利率、AI、能源、中国经济的定价均存在系统性偏差。}}
\author{KC桌面 自动生成}
\date{260606}
\begin{document}
\maketitle
\tableofcontents
\newpage
\section{一页摘要}
\begin{itemize}
  \item 美元走势出现结构性分裂：DXY指数上涨1.5\%，但贸易加权指数下跌0.1\%，背后是地缘冲突、能源贸易流和央行政策分化导致的货币分化。
  \item 美国劳动力市场韧性超预期，5月非农新增17.2万，远高于预期的8.8万，市场对美联储降息的定价可能过于乐观。
  \item 美债利率的“新地板”已经抬升，10年期收益率年末目标从4.1\%上调至4.4\%，利率分布形态而非单次加息才是关键。
  \item 中国工业生产4月增速骤降至4.1\%，约一半拖累来自石油化工行业的供给冲击，而非需求疲软。
  \item AI竞争正从算力军备转向数据资产变现，拥有专有、难复制数据库的企业（如RELX）被市场错误地归类为“AI风险股”。
  \item 全球资金流向显示，美国市场正从“避险目的地”变成“唯一增长叙事”，工业类基金成为资金最追捧的板块。
  \item 中国A股情绪降温但未崩溃，6-12个月仍有10-12\%的上行空间，支撑来自出口改善、AI资本开支和人民币升值。
  \item 集装箱航运的运价反弹是供给侧结构性再定价的结果，有效运力比名义运力紧张得多，且缺口短期内无法弥补。
  \item 日元回流叙事被高估，真正驱动汇率的不是日本邮政银行，而是GPIF的未对冲仓位，而当前利差不支持大规模回流。
  \item 中国药企在全球创新药研发链条中已占据不可替代的结构性位置，市场高估了地缘法案的影响。
\end{itemize}
\section{宏观与利率：劳动力韧性重塑降息叙事，利率“新地板”已抬升}
\textbf{市场对美联储降息的预期可能建立在一个正在被数据推翻的假设上：劳动力市场会显著放缓。5月非农就业超预期，就业增长广度改善，周期性行业表现突出，劳动力市场的韧性正在从“结构性短缺”转向“周期性复苏”。美债利率的“新地板”已经抬升，10年期收益率年末目标上调至4.4\%。}

\begin{itemize}
  \item 5月非农就业新增17.2万，远高于市场预期的8.8万，前两个月数据被上修，3个月均值创2024年3月以来新高。
  \item 就业增长的广度在改善，休闲酒店业新增7万（受世界杯临时性因素推动），地方政府就业增加5.5万，周期性行业招聘广度持续提升。
  \item 失业率稳定在4.3\%，5月失业人数下降，扭转了4月的上升趋势，但求职成功率仍偏弱。
  \item 时薪环比增长0.3\%，带动工资收入增长0.4\%，工资增长和消费支出的下行空间可能比市场预期的更有限。
  \item 核心CPI预计从4月的0.376\%降至5月的0.183\%，主要因租金技术调整消退，但核心PCE预计从0.239\%升至0.327\%，年化核心PCE预计达3.4\%，远高于美联储2\%目标。
  \item 市场对美联储政策路径的定价呈现不对称结构：加息风险前置，降息空间被压缩。5月就业报告发布后，前端利率重新定价朝一个方向倾斜。
  \item 10年期美债收益率年末目标从4.1\%上调至4.4\%，英国国债10年期目标上调至4.5\%。
  \item 欧央行预计下周加息25bp，但7月是否继续是关键。通胀远期定价近期回落，但降幅已超过宏观和能源价格隐含的水平。
  \item 日本央行行长暗示6月加息，概率升至95\%，但曲线在鹰派信号后反而陡化，说明市场对反应函数变化存疑。
  \item 加拿大央行面临两难：通胀下行但经济已衰退，短期加息门槛高。
\end{itemize}
\begin{figure}[htbp]
\centering
\includegraphics[width=0.88\linewidth]{figures/fig_004.jpg}
\caption{Exhibit 1 - Exhibit 1: The 1y1y point has typically led on the US curve during sharp inflections in front-end rates 3m change in 2y swap rate vs 1y/1y1y/2y1y swap fly return (higher = belly underperforms)}
\end{figure}
\begin{figure}[htbp]
\centering
\includegraphics[width=0.88\linewidth]{figures/fig_005.jpg}
\caption{Exhibit 2 - Exhibit 2: Solid YTD Money Fund AUM growth reflects the higher-for-longer front-end yields}
\end{figure}
\begin{figure}[htbp]
\centering
\includegraphics[width=0.88\linewidth]{figures/fig_006.jpg}
\caption{Exhibit 3 - Exhibit 3: The recent decline in traded inflation has outpaced what the macro and energy pricing imply Fair value model based on energy prices, core inflation, the euro currency and the level of nominal yields}
\end{figure}
\begin{figure}[htbp]
\centering
\includegraphics[width=0.88\linewidth]{figures/fig_007.jpg}
\caption{Exhibit 4 - Exhibit 4: Signals for Gilt spreads suggest supply is not the key issue for Gilts Fair value model based on repo-ois spread, rates volatility, free float, the level of yield and the curve}
\end{figure}
\begin{figure}[htbp]
\centering
\includegraphics[width=0.88\linewidth]{figures/fig_008.jpg}
\caption{Exhibit 5 - Exhibit 5: Hawkish BoJ policy innovations typically result in bear flattening consistent with a more front-loaded policy reaction JGB Curve Betas* to Changes in 3m OIS}
\end{figure}
{\small\color{refgray}\textbf{References:} [R002] 美国利率的真正拐点，取决于一个被忽视的变量; [R009] 就业市场比通胀更值得关注的信号：美联储的降息叙事正在被重新定价}

\section{美元与外汇：美元走势分裂，日元回流叙事被高估}
\textbf{美元不再是一个统一的“强美元”或“弱美元”故事，而是被地缘冲突、能源贸易流和各国央行政策分化撕成了两半。日元回流叙事被市场过度简化，真正可能对日元汇率产生实质性影响的不是日本邮政银行，而是GPIF的未对冲仓位。}

\begin{itemize}
  \item DXY指数（欧元权重超57\%）年初至今上涨约1.5\%，而更全面的美元贸易加权指数下跌约0.1\%，美元强弱完全取决于与谁比较。
  \item 美国国内周期性数据强劲（非农超预期、ISM调查显示增长前景改善与通胀压力并存），利率差偏向美元，尤其是相对于欧洲。
  \item 中国人民币持续升值（趋势渐进但将超出市场预期）和日元在干预威胁下的“粘性”，使美元在亚洲和新兴市场方向承压。
  \item 中东冲突导致贸易条件在不同经济体之间产生巨大分化：能源出口国货币受益，能源进口国货币承压。
  \item 若中东冲突缓和，受益者将从能源出口国切换到能源进口国：新兴市场包括南非兰特、韩元、波兰兹罗提、智利比索、匈牙利福林；G10包括瑞典克朗、新西兰元、英镑。
  \item 日本邮政银行持有约5500亿美元海外资产，但它是高度对冲的投资者，其投资决策取决于“对冲后的海外收益率”与“日债收益率”的比较，而非名义收益率。
  \item GPIF管理近1万亿美元资产，大部分未对冲汇率风险，若减持海外债券至下限（20\%），可释放约870亿美元回流日债。
  \item 对未对冲投资者来说，日美10年期利差仍为负，没有更有利的利差环境，日元升值式的资金回流很难出现。
  \item GPIF和私人养老金仍在买海外债券，最新数据未显示大规模回流迹象。
\end{itemize}
\begin{figure}[htbp]
\centering
\includegraphics[width=0.88\linewidth]{figures/fig_001.jpg}
\caption{Exhibit 1 - Exhibit 1: Divergent year-to-date performance across DXY and our Dollar TWI is testament to divided recent Dollar performance 2026 Year-to-Date GS USD TWI and DXY Composition}
\end{figure}
\begin{figure}[htbp]
\centering
\includegraphics[width=0.88\linewidth]{figures/fig_002.jpg}
\caption{Exhibit 2 - Exhibit 2: High-beta energy importer currencies would likely benefit most from a resolution to the conflict, unwinding their recent outperformance USD/XXX Vulnerabilities to Terms of Trade Shifts vs Risk-Off Shifts (Larger Bubble}
\end{figure}
\begin{figure}[htbp]
\centering
\includegraphics[width=0.88\linewidth]{figures/fig_003.jpg}
\caption{Exhibit 3 - Exhibit 3: INR carry is now higher than for other Asia high-yielders (IDR and PHP), ZAR and MXN}
\end{figure}
\begin{figure}[htbp]
\centering
\includegraphics[width=0.88\linewidth]{figures/fig_009.jpg}
\caption{Exhibit 1 - Exhibit 1: Japanese investors (excluding the MoF) hold nearly \textbackslash{}\$6tn of foreign assets, with the majority mostly unhedged}
\end{figure}
\begin{figure}[htbp]
\centering
\includegraphics[width=0.88\linewidth]{figures/fig_010.jpg}
\caption{Exhibit 2 - Data from Post Bank and GPIF financial reports, the BoJ's April FSR, and Flow of Funds data (Outward Investment in Securities + Investment Trust Beneficiary Certificates) as of Q4 2025.}
\end{figure}
\begin{figure}[htbp]
\centering
\includegraphics[width=0.88\linewidth]{figures/fig_013.jpg}
\caption{Exhibit 5 - Exhibit 5: The timeliest data suggest that GPIF and the private pensions have continued buying foreign debt while selling foreign equities}
\end{figure}
\begin{figure}[htbp]
\centering
\includegraphics[width=0.88\linewidth]{figures/fig_014.jpg}
\caption{Exhibit 6 - Exhibit 6: While JGB yields have risen more than UST yields over the past year, the differential remains negative}
\end{figure}
{\small\color{refgray}\textbf{References:} [R001] 美元真正的分裂不在指数里，而在“谁在涨、谁在跌”的背离中; [R003] 日元回流叙事被高估了：真正驱动汇率的不是日本邮政银行，而是GPIF的未对冲仓位}

\section{AI与科技：从算力军备到数据护城河，上游设备与封测受益}
\textbf{AI的竞争正在从算力军备竞赛转向数据资产的深度变现，拥有专有、难以复制的数据库的企业被市场错误地归类为“AI风险敞口”标的。中国科技板块中，上游半导体设备、封测和代工的订单可见度正在从“季度”拉长到“两年”，下游则面临成本上升和需求分化的双重挤压。}

\begin{itemize}
  \item AI的经济价值高度集中在半导体环节，但这种结构不可持续，数据架构和编排能力才是企业真正落地AI的关键。
  \item 拥有专有、难复制数据库的企业（如RELX AI框架评分9/10、Pearson 8分、Springer Nature 7.5分）被市场错误地归类为“AI风险股”。
  \item Broadcom AI半导体季度订单积压超300亿美元，订单能见度延伸至2028年，客户提前下单不仅为锁定晶圆和HBM产能，更为确保电力基础设施配套。
  \item AI基础设施的供给约束正在从“芯片产能”向“数据中心电力容量”迁移，3至5年的长期协议已覆盖基板、HBM和CCL等关键瓶颈环节。
  \item 中国科技板块中，半导体子行业优先顺序重新排列为：半导体设备 > 封测 > 代工 > 无晶圆设计。
  \item 2026年，本土WFE供应商来自存储和先进逻辑的需求预计分别同比增长超50\%和30\%，2027年还有上行空间。
  \item 长鑫存储和长江存储的IPO时间表推进，使资本开支计划从“不确定”变为“可预期”，订单可见度拉长。
  \item 国产AI芯片从2026年下半年进入量产爬坡期，将带动后端设备（先进封装、测试等）的增量需求。
  \item 封测/代工环节产能利用率高，2Q26持续涨价，AI相关（GPU、PMIC）的封测厂受益更大。
  \item 下游消费电子（安卓手机、AIoT）需求走弱，依赖消费电子又无法转嫁成本的Fabless公司压力最大。
\end{itemize}
\begin{figure}[htbp]
\centering
\includegraphics[width=0.88\linewidth]{figures/fig_122.jpg}
\caption{Exhibit 1 - Exhibit 1: Overview of Adam Berlin's AI moat scores by sector across stocks that have shown a correlation with the GS AI basket GS structural AI moat score (1-10) - Higher score is a higher moat}
\end{figure}
\begin{figure}[htbp]
\centering
\includegraphics[width=0.88\linewidth]{figures/fig_123.jpg}
\caption{Exhibit 18 - Louisa Qiu \# Specialist Sales Christian Moore This week's scripts were (as expected, and flagged previously) impacted by the long weekend, and hence total scripts were down \textbackslash{}\textasciitilde{}6\% for the class. Foundayo saw 16,982 TRx in its eighth week, +971 and +6.1\% sequenti}
\end{figure}
{\small\color{refgray}\textbf{References:} [R015] 市场真正低估的不是AI需求，而是数据护城河的再定价; [R016] 中国科技股的核心机会不在下游，而在上游的设备与封测; [R019] AI基础设施的供给瓶颈，正在从芯片设计延伸至整个物理供应链}

\section{能源与大宗：供给侧再定价驱动运价与锂价，中国油价调控成本可控}
\textbf{集装箱航运运价反弹的核心驱动力是供给侧的结构性再定价，而非需求意外强劲。锂价短期回调不是反转，CATL新产能投产将带来可量化的供需错配窗口。中国油价调控的财政成本年化不到GDP的0.3\%，且被新能源汽车渗透加速稀释。}

\begin{itemize}
  \item SCFI指数自3月以来飙升73\%，跨太平洋和亚欧航线录得两位数周涨幅，但核心驱动力是供给侧的结构性再定价。
  \item 有效运力比名义运力紧张得多：港口拥堵、航线绕行（红海绕行好望角）及船公司主动空白航行管理吸收大量可用运力。
  \item Clarksons港口拥堵指数在2026年1-5月创历史新高，全球相当比例集装箱船队处于“等待”或“绕行”状态。
  \item 亚洲船公司有额外红利：高硫/低硫燃料价差扩大至135美元/吨，长荣脱硫塔安装率95\% vs 行业47\%，燃油成本传导能力更强。
  \item 碳酸锂价格从5月底17.55万元/吨回落至16.83万元/吨，但基本面并未松动，库存结构改善：下游正极材料厂库存环比+7\%，冶炼厂和贸易商库存分别-2\%和-8\%。
  \item CATL新产能将在2026年8-9月前投产，届时将带来单月约5000-6000吨的增量锂需求，供需只会更紧。
  \item 中国成品油定价机制以国际原油价格每桶40-130美元为边界的非线性传导系统，具备“自动稳定器”功能。
  \item 历史数据看，布伦特涨10\%，国内汽油柴油约涨5\%。中东冲突后3-4月国内油品零售量同比降17\%，价格涨17\%。
  \item 政府干预成本年化约0.3\%GDP，库存增值和上游利润能对冲，成本可控。新能源车渗透率已到54\%，替代效应在加速。
  \item 中国进口多元化：中东进口降38\%，俄罗斯+13\%，拉美+64\%，总量只降11\%。
\end{itemize}
\begin{figure}[htbp]
\centering
\includegraphics[width=0.88\linewidth]{figures/fig_064.jpg}
\caption{Figure 1 - Figure 1: Global Liners' YTD Share Price performance (2026)}
\end{figure}
\begin{figure}[htbp]
\centering
\includegraphics[width=0.88\linewidth]{figures/fig_067.jpg}
\caption{Figure 4 - Figure 4: \% of Capacity Scrubber-fitted Across Industry (Based on Clarksons' Vessel Ownership Tracking)}
\end{figure}
\begin{figure}[htbp]
\centering
\includegraphics[width=0.88\linewidth]{figures/fig_068.jpg}
\caption{Figure 5 - Figure 5: HFSO vs VLSFO Bunker Fuel Price (USD/ton)}
\end{figure}
\begin{figure}[htbp]
\centering
\includegraphics[width=0.88\linewidth]{figures/fig_069.jpg}
\caption{Figure 6 - Figure 6: Shanghai Containerized Freight Index by route}
\end{figure}
\begin{figure}[htbp]
\centering
\includegraphics[width=0.88\linewidth]{figures/fig_079.jpg}
\caption{Figure 16 - Figure 16: Clarksons Port Congestion Index - Containerships in Port, m.TEU, 7dma Figure 17: Port Congestion Index - Containerships in Port, China P.R., m.TEU, 7dma Figure 18: Port Congestion Index - Containerships in Port, Sout}
\end{figure}
\begin{figure}[htbp]
\centering
\includegraphics[width=0.88\linewidth]{figures/fig_094.jpg}
\caption{Exhibit 2 - Exhibit 1: Domestic retail fuel prices have risen less than Brent crude oil prices since the start of the Middle East conflict}
\end{figure}
\begin{figure}[htbp]
\centering
\includegraphics[width=0.88\linewidth]{figures/fig_095.jpg}
\caption{Exhibit 2 - Exhibit 2: Increased imports from Russia and Latin America have partially offset falling Middle East supply}
\end{figure}
\begin{figure}[htbp]
\centering
\includegraphics[width=0.88\linewidth]{figures/fig_096.jpg}
\caption{Exhibit 3 - Exhibit 3: An illustration of China's domestic retail oil product pricing framework ```mermaid graph TD}
\end{figure}
\begin{figure}[htbp]
\centering
\includegraphics[width=0.88\linewidth]{figures/fig_099.jpg}
\caption{Exhibit 8 - Exhibit 8: For every \$10\textbackslash{}\%\$ domestic retail oil price increase, China's domestic oil product retail sales volume would decline by around \$5\textbackslash{}\%\$}
\end{figure}
\begin{figure}[htbp]
\centering
\includegraphics[width=0.88\linewidth]{figures/fig_100.jpg}
\caption{Exhibit 9 - Exhibit 9: Domestic retail oil product sales volume contracted by 17\% yoy during March-April whereas retail prices gained 17\%}
\end{figure}
\begin{figure}[htbp]
\centering
\includegraphics[width=0.88\linewidth]{figures/fig_105.jpg}
\caption{Figure 1 - \#\# Shreyas Madabushi Figure 1. Lithium carbonate monthly inventory}
\end{figure}
\begin{figure}[htbp]
\centering
\includegraphics[width=0.88\linewidth]{figures/fig_106.jpg}
\caption{Figure 2 - © 2026 Citi Inc. No redistribution without Citi's written permission. Figure 2. Lithium carbonate weekly inventory}
\end{figure}
{\small\color{refgray}\textbf{References:} [R007] 市场真正低估的不是需求，而是供给侧的再定价; [R011] 中国油价调控的真正成本，比市场担心的低得多; [R012] 锂价短期回调不是反转，真正被低估的是CATL产能释放带来的需求硬缺口}

\section{中国经济：供给冲击而非需求疲软，消费回暖但地产仍弱}
\textbf{4月工业增加值增速骤降至4.1\%，约一半拖累来自石油化工行业的供给冲击（原材料断供），而非需求疲软。消费端高频指标改善，但地产销售和工业活动仍偏弱，政策发力节奏需要关注。}

\begin{itemize}
  \item 4月工业增加值同比增速从5.7\%骤降至4.1\%，远低于市场预期的6.0\%，创近三年最低月度纪录。
  \item 石油、煤炭及其他燃料加工业产出增速从3月的8.1\%骤降至-0.9\%，原油加工量同比收缩5.8\%，沥青产出量同比暴跌40.1\%。
  \item 化工原材料及制品业产出增速从9.0\%放缓至5.3\%，化学纤维制造业产出增速降至2.2\%，为2023年3月以来最低。
  \item 综合估算，油化板块直接拖累工业增速0.7个百分点，加上间接影响，贡献了4月工业减速的一半左右。
  \item AI相关电子设备增速逆势上行至15.6\%，拉动0.4个百分点，如果没有这块，工业增速会更难堪。
  \item 高频数据暗示5月供应瓶颈仍在延续：山东地炼开工率从60\%降至54.7\%，沥青工厂开工率跌至16\%，PTA开工率从69\%骤降至59.3\%。
  \item 消费端数据改善：Morning Consult消费者信心指数升至多年新高，30城新房成交面积高于去年同期，二手房成交面积维持高位。
  \item 国内汽油和柴油价格6月4日分别下调525元/吨和505元/吨，降低居民出行成本。
  \item 地方专项债累计发行约1500亿元，进度偏慢；PSL余额5月继续收缩（净偿还）。
  \item 投行预计5月工业增加值同比增速维持在4\%左右低位，Q2 GDP增速预测下调至4.1\%（市场共识4.7\%）。
\end{itemize}
\begin{figure}[htbp]
\centering
\includegraphics[width=0.88\linewidth]{figures/fig_015.jpg}
\caption{Exhibit 1 - Exhibit 1: 30-city daily property transaction volume in the primary market was roughly flat over the last week but remained above year-ago level}
\end{figure}
\begin{figure}[htbp]
\centering
\includegraphics[width=0.88\linewidth]{figures/fig_016.jpg}
\caption{Exhibit 2 - Exhibit 2: 16-city daily property transaction volume in the secondary market declined but remained above year-ago level}
\end{figure}
\begin{figure}[htbp]
\centering
\includegraphics[width=0.88\linewidth]{figures/fig_022.jpg}
\caption{Exhibit 8 - Exhibit 8: The Morning Consult consumer confidence edged up to multi-year high}
\end{figure}
\begin{figure}[htbp]
\centering
\includegraphics[width=0.88\linewidth]{figures/fig_024.jpg}
\caption{Exhibit 10 - Exhibit 10: Steel demand edged down over the past week}
\end{figure}
\begin{figure}[htbp]
\centering
\includegraphics[width=0.88\linewidth]{figures/fig_025.jpg}
\caption{Exhibit 11 - Exhibit 11: Steel production edged down over the past week}
\end{figure}
\begin{figure}[htbp]
\centering
\includegraphics[width=0.88\linewidth]{figures/fig_026.jpg}
\caption{Exhibit 12 - Exhibit 12: Daily coal consumption in coastal provinces increased further over the last week and remained above year-ago level}
\end{figure}
\begin{figure}[htbp]
\centering
\includegraphics[width=0.88\linewidth]{figures/fig_027.jpg}
\caption{Exhibit 13 - Exhibit 13: RMB 1.50bn local government special bonds have been issued year-to-date}
\end{figure}
\begin{figure}[htbp]
\centering
\includegraphics[width=0.88\linewidth]{figures/fig_028.jpg}
\caption{Exhibit 14 - Exhibit 14: PSL loans outstanding contracted in May 2026 PSL: Net injection vs outstanding amount}
\end{figure}
\begin{figure}[htbp]
\centering
\includegraphics[width=0.88\linewidth]{figures/fig_111.jpg}
\caption{Figure 5 - - The operating rate of purified terephthalic acid (PTA) factories fell sharply to 59.3\% at end-May from 69.0\% at end-April (Figure 5), before rebounding moderately to 66.4\% in early June. The monthly average operating rate dropped further to 8.9pp below last }
\end{figure}
\begin{figure}[htbp]
\centering
\includegraphics[width=0.88\linewidth]{figures/fig_112.jpg}
\caption{Figure 7 - Before the Iran war, \$50\textbackslash{}\%\$ of China's oil imports and \$16\textbackslash{}\%\$ of its natural gas imports passed through the Strait of Hormuz (SoH). Due to the severe damage to a wide range of energy facilities in the Middle East and the blockage of the SoH, the supply disrupt}
\end{figure}
\begin{figure}[htbp]
\centering
\includegraphics[width=0.88\linewidth]{figures/fig_113.jpg}
\caption{Figure 9 - \#\# The import disruptions of chemical raw materials For sulfur, import volumes crashed by \$72.4\textbackslash{}\%\$ y-o-y in April to its lowest monthly reading since October 2008 (Figure 9), and the supply shortage worsened even further after April. Inventories of sulfur at C}
\end{figure}
\begin{figure}[htbp]
\centering
\includegraphics[width=0.88\linewidth]{figures/fig_114.jpg}
\caption{Figure 11 - On essential chemical raw materials for producing plastics and fabrics, import volumes of ethylene glycol plummeted by \$46.3\textbackslash{}\%\$ y-o-y in April to its lowest monthly reading since 2007 (Figure 11). Import volumes of methanol dropped by \$28.6\textbackslash{}\%\$ y-o-y in April a}
\end{figure}
{\small\color{refgray}\textbf{References:} [R004] 中国经济的真实温度：需求在修复，但供给侧的再定价才是关键; [R013] 市场真正低估的不是需求，而是供给侧的再定价}

\section{全球资金流向：美国成唯一增长叙事，工业板块受追捧}
\textbf{全球资金正在从“防御性集中”转向“结构性再配置”，美国市场从“避险目的地”变成“唯一增长叙事”。工业类基金成为资金最追捧的板块，而金融和消费品板块遭遇最大规模流出。}

\begin{itemize}
  \item 截至6月3日当周，全球股票基金净流入230亿美元，扭转前一周70亿美元的净流出，几乎全部由美国市场贡献。
  \item 除美国外，发达市场普遍呈现净流出，资金不是在“抄底全球”，而是在“加注美国”。
  \item 美国股票基金净流入在经历4月剧烈波动后，4周移动均值已回到正值区间，修复速度比市场预期更快。
  \item 工业类基金成为资金最追捧的板块，累计净流入占管理资产比例高达24\%。
  \item 金融和消费品板块遭遇最大规模流出，科技板块表现中性，累计流入约3\%。
  \item 全球固定收益基金当周净流入超400亿美元，但内部结构变化：短期债券和通胀保护债券持续受欢迎，长久期债券净流出。
  \item 货币市场基金资产单周飙升1220亿美元，说明市场仍有大量资金在观望。
  \item 外汇流向方面，美元需求最强，人民币净流出最多，英镑有亮点。
  \item 新兴市场内部资金分化明显：全球新兴市场基准基金、中国A股基金、韩国基金流出，只有台湾基金逆势流入。
\end{itemize}
{\small\color{refgray}\textbf{References:} [R005] 资金流向正在告诉你一个被忽视的结构性切换}

\section{A股与港股：情绪降温但未崩溃，结构性信号尚未被定价}
\textbf{A股情绪指标MSASI从66\%回落至62\%，但情绪下滑来自“参与意愿的边际递减”而非恐慌性抛售。6-12个月仍有10-12\%的上行空间，支撑来自出口改善、AI/能源资本开支、人民币升值及互联互通机制优化。}

\begin{itemize}
  \item MSASI加权值从66\%降至62\%，30日RSI仅上升1\%，融资余额基本未变，情绪下滑来自参与意愿的边际递减。
  \item 创业板日均成交额降至7790亿元，全市场降至3.02万亿元，环比分别下降9\%和6\%，但股指期货成交额和融资余额未同步萎缩。
  \item 5月制造业PMI滑至50\%荣枯线，消费出口订单和能源密集型生产走弱，但高端制造依然坚挺，K型分化格局持续。
  \item PPI环比回落至1\%（4月为1.7\%），油价涨势放缓。国内经济团队认为4-5月数据偏弱，Q2 GDP下行压力加大，预计6月起财政加速落地。
  \item 南向资金持续流入：5月28日-6月3日净流入45亿美元，年内累计353亿美元。
  \item 港股医药板块在ASCO数据公布后呈现“卖事实”弱势，但中国药企在全球创新药研发链条中已占据不可替代的结构性位置。
  \item 2025年中国医药研发支出预计达390亿美元，占全球约13\%，自2020年起每年进入临床的创新药数量已全球第一。
  \item 2026年前5个月，中国license-out首付款达48亿美元，同比增长93\%，5月单月首付款12.5亿美元。
  \item 若美国切断与中国药企的合作，受损更大的可能是美国自身：2020-2026年中国license-out交易中，美国企业占比42\%但贡献了57\%的首付款。
\end{itemize}
\begin{figure}[htbp]
\centering
\includegraphics[width=0.88\linewidth]{figures/fig_036.jpg}
\caption{Exhibit 1 - Exhibit 1: MS A-share Sentiment Indicator: MSASI weighted and MSASI weighted 1MMA}
\end{figure}
\begin{figure}[htbp]
\centering
\includegraphics[width=0.88\linewidth]{figures/fig_039.jpg}
\caption{Exhibit 4 - Exhibit 4: ChiNext turnover adjusted by moving 100D min-max (scaled to 0-100\% based on the percentage away from its 100-day high and low levels) vs. CSI 300 relative to 100D MA}
\end{figure}
\begin{figure}[htbp]
\centering
\includegraphics[width=0.88\linewidth]{figures/fig_040.jpg}
\caption{Exhibit 5 - Exhibit 5: A-share turnover adjusted by moving 100D min-max (scaled to 0-100\% based on the percentage away from its 100-day high and low levels) vs. CSI 300 relative to 100D MA}
\end{figure}
\begin{figure}[htbp]
\centering
\includegraphics[width=0.88\linewidth]{figures/fig_115.jpg}
\caption{Figure 1 - Figure 1: YTD performance of HK-listed healthcare companies}
\end{figure}
\begin{figure}[htbp]
\centering
\includegraphics[width=0.88\linewidth]{figures/fig_116.jpg}
\caption{Figure 2 - Figure 2: Global and China R\&D spending (USD bn)}
\end{figure}
\begin{figure}[htbp]
\centering
\includegraphics[width=0.88\linewidth]{figures/fig_117.jpg}
\caption{Figure 3 - Figure 3: Number of innovative drugs newly entered into clinical trials}
\end{figure}
\begin{figure}[htbp]
\centering
\includegraphics[width=0.88\linewidth]{figures/fig_118.jpg}
\caption{Figure 4 - Figure 4: Geographical breakdown of out-licensing deals Geographical breakdown of out-licensing deals}
\end{figure}
\begin{figure}[htbp]
\centering
\includegraphics[width=0.88\linewidth]{figures/fig_119.jpg}
\caption{Figure 5 - Figure 5: Geographical breakdown of upfront payments Geographical breakdown of upfront}
\end{figure}
\begin{figure}[htbp]
\centering
\includegraphics[width=0.88\linewidth]{figures/fig_120.jpg}
\caption{Figure 6 - Figure 6: Out-licensing deals in China}
\end{figure}
{\small\color{refgray}\textbf{References:} [R006] A股情绪降温背后，真正的结构性信号尚未被定价; [R014] 市场真正低估的不是地缘风险，而是中国药企已不可替代的全球研发节点地位}

\section{欧洲汽车与日本IT：结构性变化下的赢家与输家}
\textbf{中国车企在欧洲市场份额翻倍至近7\%，但驱动力是产品力而非价格战。欧洲车企选择保护利润而非份额，这正在重塑竞争逻辑。日本汽车行业IT投资正在快速降温，但金融行业核心系统更新需求成为替代增长引擎。}

\begin{itemize}
  \item 中国车企在欧洲市场份额从2023年初约1\%攀升至2026年4月近7\%，增速远超历史上任何新进入者。
  \item 在英国市场，奇瑞和比亚迪在极短时间内达到3-4\%市场份额，而丰田和现代当年达到这一水平耗费数倍时间。
  \item 中国车企竞争力并非来自价格战：在德国市场，中国品牌定价基本处于各细分市场的中位区间，驱动消费者选择的是产品“物超所值”。
  \item 反补贴关税未挡住中国车企：中国车企直接消化了关税成本，欧洲市场60\%是公司车队采购，残值重要，中国品牌未疯狂压价。
  \item 欧洲车企选择保护利润而非份额：Stellantis、福特、日韩品牌份额流失最严重，但欧洲车企更倾向保护现金流和利润率，加速用便宜的LFP电池降成本。
  \item 日本汽车行业IT投资正在经历周期性回调：Systena称部分整车厂投资意愿减弱，Argo Graphics六年来首次营业利润同比下降，DIT嵌入式解决方案业务销售额和毛利润均同比下降。
  \item 日本金融行业核心系统更新需求成为替代增长引擎，那些服务多行业的系统集成商（如富士通、NEC）值得关注。
\end{itemize}
{\small\color{refgray}\textbf{References:} [R010] 欧洲汽车市场真正被低估的风险，不是中国份额，而是欧洲的利润优先策略; [R023] 日本汽车业IT投资正在踩刹车，但这恰恰暴露了另一条增长主线}

\section{医药与GLP-1：口服药渗透率加速，多肽制造绿色化成瓶颈}
\textbf{GLP-1口服药的处方数据捕获率差异正在扭曲市场对两款核心药物相对表现的认知。多肽API制造的高溶剂消耗和高过程质量强度，正在成为制约产能有效扩张的隐形天花板。}

\begin{itemize}
  \item 截至5月29日当周，GLP-1类药物总处方量244万张，环比下降约6\%（长周末季节性扰动），4周同比增速仍维持在39.2\%高位。
  \item 礼来整体（Zepbound、Mounjaro、Foundayo合计）在总处方量中占比59.7\%，诺和诺德为40.3\%。
  \item Foundayo第8周处方量16,982，环比增长6.1\%，仍在爬坡期，但IQVIA捕获率仅约70\%。
  \item 诺和诺德在股东大会上透露：Wegovy口服第10周处方量是Zepbound的3倍，但IQVIA数据仅显示1.6倍，IQVIA对Wegovy口服的捕获率可能仅50\%。
  \item 到2030年，多肽药物市场规模预计达1200亿美元，年复合增长率18\%，GLP-1类药物是核心驱动力。
  \item 生产1公斤GLP-1类多肽API平均消耗约14吨溶剂，而典型小分子药物仅0.3吨，多肽API被标记为“最不可持续的药物模态之一”。
  \item 头部CDMO正通过连续色谱技术（溶剂消耗降30\%，材料使用降70\%）、溶剂回收和绿色合成路线降本增效。
  \item ASCO 2026上，PD-1/VEGF双抗竞争进入“数据细节决定估值分化”阶段，IL-2赛道正在经历“可成药性基准”的转移。
\end{itemize}
\begin{figure}[htbp]
\centering
\includegraphics[width=0.88\linewidth]{figures/fig_124.jpg}
\caption{Exhibit 8 - EXHIBIT 1: Oral GLP-1 Launch Trajectories (TRx)}
\end{figure}
\begin{figure}[htbp]
\centering
\includegraphics[width=0.88\linewidth]{figures/fig_125.jpg}
\caption{EXHIBIT 2 - EXHIBIT 2: Wegovy Pill \& Foundayo vs Zepbound Launch (New to Brand - NBRx; Weekly)}
\end{figure}
\begin{figure}[htbp]
\centering
\includegraphics[width=0.88\linewidth]{figures/fig_128.jpg}
\caption{EXHIBIT 7 - EXHIBIT 7: We estimate \textbackslash{}\textasciitilde{}70\% capture rate for Foundayo EXHIBIT 8: Early data points, but daily Foundayo scripts - although materially lower - look to trend well with weekly scripts}
\end{figure}
\begin{figure}[htbp]
\centering
\includegraphics[width=0.88\linewidth]{figures/fig_130.jpg}
\caption{EXHIBIT 11 - EXHIBIT 11: At the 2026 Novo AGM, Novo reported that weekly Wegovy pill TRx were three times higher than tirzepatide at week 10}
\end{figure}
\begin{figure}[htbp]
\centering
\includegraphics[width=0.88\linewidth]{figures/fig_150.jpg}
\caption{EXHIBIT 1 - EXHIBIT 1: Oncology and obesity are the fastest growing therapeutic areas - and should remain so in the near term Global historic and forecast growth for top 20 therapy areas}
\end{figure}
\begin{figure}[htbp]
\centering
\includegraphics[width=0.88\linewidth]{figures/fig_151.jpg}
\caption{EXHIBIT 2 - EXHIBIT 2: Depending on pricing and reimbursement developments, obesity spending could break into the \textbackslash{}\$100Bn level in the coming years Global obesity spending and growth}
\end{figure}
\begin{figure}[htbp]
\centering
\includegraphics[width=0.88\linewidth]{figures/fig_152.jpg}
\caption{EXHIBIT 3 - EXHIBIT 3: Among the newer molecular modalities, peptides stand out as both the largest market by 2030 and one of the fastest-growing Market size forecast by modality}
\end{figure}
\begin{figure}[htbp]
\centering
\includegraphics[width=0.88\linewidth]{figures/fig_153.jpg}
\caption{EXHIBIT 5 - EXHIBIT 5: Comparison of peptide synthesis methods - despite high solvent waste, SPPS is the commercial standard for mid to long peptides like GLP-1 agonists EXHIBIT 6: Major CDMOs are expanding their peptide synthesis capacity to}
\end{figure}
\begin{figure}[htbp]
\centering
\includegraphics[width=0.88\linewidth]{figures/fig_154.jpg}
\caption{EXHIBIT 8 - EXHIBIT 8: Compared to traditional batch chromatography, continuous chromatography (multi-column countercurrent solvent gradient purification or MCSGP) increases throughput and reduces process control and lyphilization volume Thr}
\end{figure}
{\small\color{refgray}\textbf{References:} [R017] GLP-1口服药的战场，IQVIA数据正在低估真实竞争烈度; [R020] ASCO 2026第三天：市场低估的不是PD-1/VEGF组合的竞争，而是IL-2赛道“可成药性基准”的转移; [R021] 减肥药热潮的真正瓶颈不在需求，而在多肽制造的绿色化能力}

\section{保险与支付：渠道重塑与信任定价}
\textbf{中国保险行业基本面确实在变好，但变好的方式比多数人想象的更微妙——银保渠道费用监管收紧正在悄悄改变竞争格局，让大型保险公司重新拿回定价权。支付行业的核心不是费率高低，而是“信任即服务”的定价权。}

\begin{itemize}
  \item 2025年3月银保渠道费用新规出台，大型保险公司普遍欢迎：更严格的费用审查推动银行倾向与大保司合作。
  \item 中国太平银保渠道VONB贡献约35\%，管理层明确表示有潜力逐步提升到50\%。
  \item 国寿经历2026年一季度VONB同比增长75\%的异常高增长后，管理层预计后续增速回归正常。
  \item 保险公司投资策略以红利股为主线：目标股息率普遍在4\%以上，部分公司OCI账户中权益投资占比已升至约30\%。
  \item 财险承保空间仍在，权益市场反弹带动投资收益和净资产恢复。
  \item 美国信用卡平均MDR为2.4\%，借记卡仅0.7\%，但信用卡对应的是集成了信用、奖励、争议处理、欺诈保障的复合服务产品。
  \item 信用卡2.4\%费率中，真正落到卡组织手里的只有0.23\%，大头流向发卡行覆盖奖励积分（约3\%）、欺诈损失（约1.5\%）、争议处理（约1.2\%）、资金成本（约2.5\%）。
  \item 欧洲2015年将借记卡费率限制在0.2\%、信用卡0.3\%后，卡渗透率从33\%飙升至77\%，说明合理费率上限能刺激更多商户接受卡片。
  \item 美国34\%的小商户开始对信用卡交易额外收费（surcharging），但令牌化（tokenization）正在巩固卡组织的定价权。
\end{itemize}
\begin{figure}[htbp]
\centering
\includegraphics[width=0.88\linewidth]{figures/fig_155.jpg}
\caption{Exhibit 1 - EXHIBIT 1: Debit is significantly cheaper vs. credit in the U.S. Weighted Average Merchant Card Fees in the US (2025, \%)}
\end{figure}
\begin{figure}[htbp]
\centering
\includegraphics[width=0.88\linewidth]{figures/fig_156.jpg}
\caption{Exhibit 2 - EXHIBIT 2: Cards are especially cheaper in Europe Merchant Discount Rate by Country}
\end{figure}
\begin{figure}[htbp]
\centering
\includegraphics[width=0.88\linewidth]{figures/fig_157.jpg}
\caption{EXHIBIT 3 - EXHIBIT 3: Mix between credit/debit cards varies widely by country.. Debit vs Credit Mix in Select Countries (\%, 2025)}
\end{figure}
\begin{figure}[htbp]
\centering
\includegraphics[width=0.88\linewidth]{figures/fig_161.jpg}
\caption{Exhibit 7 - EXHIBIT 7: Europe interchange regulation went into effect in Dec 2015, significantly boosting card penetration and adoption Europe Card Penetration Rate}
\end{figure}
\begin{figure}[htbp]
\centering
\includegraphics[width=0.88\linewidth]{figures/fig_162.jpg}
\caption{Exhibit 8 - EXHIBIT 8: Surcharging among small business owners reached 34\% per a survey conducted early 2025 \% Small Business Owners Adopting Surcharging}
\end{figure}
\begin{figure}[htbp]
\centering
\includegraphics[width=0.88\linewidth]{figures/fig_163.jpg}
\caption{Exhibit 9 - EXHIBIT 9: The networks yields are only a small fraction of total MDR US Merchant Discount Rate vs Visa US Yield}
\end{figure}
\begin{figure}[htbp]
\centering
\includegraphics[width=0.88\linewidth]{figures/fig_164.jpg}
\caption{EXHIBIT 10 - EXHIBIT 10: Tokens are proliferating, and will do so even more with Agentic commerce - solidifying moat and pricing power Visa \# of Tokens and Cards (B, CY) '20-25 CAGR Tokens: 77\% Cards: 7\%}
\end{figure}
{\small\color{refgray}\textbf{References:} [R018] 保险股反弹背后，市场真正低估的是渠道结构重塑的定价权; [R022] 市场真正低估的不是卡组织费率，而是“信任即服务”的定价权}

\section{AI重塑全球收入分配：高收入经济体才是最大输家}
\textbf{当生成式AI从“生产力工具”进化为“认知替代者”，高收入经济体才是AI冲击下最大的输家，而市场当前的定价恰恰呈现了“认知倒挂”。}

\begin{itemize}
  \item 高收入国家中，信息产业直接雇佣的劳动力占总量7\%-10\%，而中等收入国家仅1\%-4\%，发达经济体对AI的暴露度更高。
  \item 软件工程师年薪：印度1万美元 vs 美国13.6万 vs 英国8.1万，企业用AI替代高薪岗位的经济诱因远大于替代低薪岗位。
  \item 中低收入国家国民收入平均损失10\%，高收入国家国民收入平均损失22\%，AI反而可能成为“拉平”全球收入差距的力量。
  \item 印度、印尼股市在AI浪潮中跌得最惨，原因是这些市场没有AI产业链上的受益者（不造芯片、不搞模型），市场只看到短期冲击。
  \item 信息产业就业占比决定了谁真正暴露在AI冲击之下：一个经济体只有3\%的人在信息产业，AI替代一半影响1.5个百分点；若比例是10\%，替代一半就是5个百分点。
\end{itemize}
\begin{figure}[htbp]
\centering
\includegraphics[width=0.88\linewidth]{figures/fig_085.jpg}
\caption{EXHIBIT 1 - EXHIBIT 1: Middle income economies have far lower workforce share in information industries compared to high income ones Employment Share (\%) in Information Industries}
\end{figure}
\begin{figure}[htbp]
\centering
\includegraphics[width=0.88\linewidth]{figures/fig_086.jpg}
\caption{EXHIBIT 2 - EXHIBIT 2: Not only employment, but in overall share of GDP too, emerging economies are less reliant on services than high income ones}
\end{figure}
\begin{figure}[htbp]
\centering
\includegraphics[width=0.88\linewidth]{figures/fig_087.jpg}
\caption{EXHIBIT 3 - EXHIBIT 3: Emerging economies have salaries an order of magnitude below those of high income economies for the same roles Average Salaries across exposed job categories in key economies}
\end{figure}
\begin{figure}[htbp]
\centering
\includegraphics[width=0.88\linewidth]{figures/fig_088.jpg}
\caption{EXHIBIT 4 - EXHIBIT 4: Markets of middle income economies (highlighted in red) have seen the steepest falls since the spurt of AI announcements USD Returns of Indices since Anthropic's Cowork Tools Announcement}
\end{figure}
\begin{figure}[htbp]
\centering
\includegraphics[width=0.88\linewidth]{figures/fig_089.jpg}
\caption{EXHIBIT 5 - EXHIBIT 5: India has among the highest exposure to IT services in its benchmark indices Benchmark Indices exposure to IT services across key economies}
\end{figure}
\begin{figure}[htbp]
\centering
\includegraphics[width=0.88\linewidth]{figures/fig_090.jpg}
\caption{Exhibit 6 - EXHIBIT 6: Limitless productivity increase does not increase demand infinitely; it rather gradually kills it}
\end{figure}
\begin{figure}[htbp]
\centering
\includegraphics[width=0.88\linewidth]{figures/fig_091.jpg}
\caption{Exhibit 7 - EXHIBIT 7: When it comes to income inequality, middle income economies are far less equal - with top 10\% having over 50\% of national income. The corresponding number for high income economies is 43\%}
\end{figure}
\begin{figure}[htbp]
\centering
\includegraphics[width=0.88\linewidth]{figures/fig_092.jpg}
\caption{EXHIBIT 8 - EXHIBIT 8: However, when it comes to wealth, the countries are relatively closely bundled, with 60\%+ share of top 10 for all}
\end{figure}
\begin{figure}[htbp]
\centering
\includegraphics[width=0.88\linewidth]{figures/fig_093.jpg}
\caption{EXHIBIT 9 - EXHIBIT 9: IT stocks sit at their cheapest valuations since COVID period six years back}
\end{figure}
{\small\color{refgray}\textbf{References:} [R008] 市场真正低估的不是AI需求，而是它如何重塑全球收入分配}

\section{结语}
综合来看，当前市场定价存在多重结构性偏差：美元走势分裂、利率新地板抬升、AI竞争从算力转向数据、中国供给冲击而非需求疲软、以及全球资金向美国集中。这些偏差的修复过程，将定义未来6-12个月的主要资产定价逻辑。
\newpage
\section{报告覆盖清单}
\begin{itemize}
  \item [R001] 美元与外汇：美元走势分裂，日元回流叙事被高估 - 美元真正的分裂不在指数里，而在“谁在涨、谁在跌”的背离中
  \item [R002] 宏观与利率：劳动力韧性重塑降息叙事，利率“新地板”已抬升 - 美国利率的真正拐点，取决于一个被忽视的变量
  \item [R003] 美元与外汇：美元走势分裂，日元回流叙事被高估 - 日元回流叙事被高估了：真正驱动汇率的不是日本邮政银行，而是GPIF的未对冲仓位
  \item [R004] 中国经济：供给冲击而非需求疲软，消费回暖但地产仍弱 - 中国经济的真实温度：需求在修复，但供给侧的再定价才是关键
  \item [R005] 全球资金流向：美国成唯一增长叙事，工业板块受追捧 - 资金流向正在告诉你一个被忽视的结构性切换
  \item [R006] A股与港股：情绪降温但未崩溃，结构性信号尚未被定价 - A股情绪降温背后，真正的结构性信号尚未被定价
  \item [R007] 能源与大宗：供给侧再定价驱动运价与锂价，中国油价调控成本可控 - 市场真正低估的不是需求，而是供给侧的再定价
  \item [R008] AI重塑全球收入分配：高收入经济体才是最大输家 - 市场真正低估的不是AI需求，而是它如何重塑全球收入分配
  \item [R009] 宏观与利率：劳动力韧性重塑降息叙事，利率“新地板”已抬升 - 就业市场比通胀更值得关注的信号：美联储的降息叙事正在被重新定价
  \item [R010] 欧洲汽车与日本IT：结构性变化下的赢家与输家 - 欧洲汽车市场真正被低估的风险，不是中国份额，而是欧洲的利润优先策略
  \item [R011] 能源与大宗：供给侧再定价驱动运价与锂价，中国油价调控成本可控 - 中国油价调控的真正成本，比市场担心的低得多
  \item [R012] 能源与大宗：供给侧再定价驱动运价与锂价，中国油价调控成本可控 - 锂价短期回调不是反转，真正被低估的是CATL产能释放带来的需求硬缺口
  \item [R013] 中国经济：供给冲击而非需求疲软，消费回暖但地产仍弱 - 市场真正低估的不是需求，而是供给侧的再定价
  \item [R014] A股与港股：情绪降温但未崩溃，结构性信号尚未被定价 - 市场真正低估的不是地缘风险，而是中国药企已不可替代的全球研发节点地位
  \item [R015] AI与科技：从算力军备到数据护城河，上游设备与封测受益 - 市场真正低估的不是AI需求，而是数据护城河的再定价
  \item [R016] AI与科技：从算力军备到数据护城河，上游设备与封测受益 - 中国科技股的核心机会不在下游，而在上游的设备与封测
  \item [R017] 医药与GLP-1：口服药渗透率加速，多肽制造绿色化成瓶颈 - GLP-1口服药的战场，IQVIA数据正在低估真实竞争烈度
  \item [R018] 保险与支付：渠道重塑与信任定价 - 保险股反弹背后，市场真正低估的是渠道结构重塑的定价权
  \item [R019] AI与科技：从算力军备到数据护城河，上游设备与封测受益 - AI基础设施的供给瓶颈，正在从芯片设计延伸至整个物理供应链
  \item [R020] 医药与GLP-1：口服药渗透率加速，多肽制造绿色化成瓶颈 - ASCO 2026第三天：市场低估的不是PD-1/VEGF组合的竞争，而是IL-2赛道“可成药性基准”的转移
  \item [R021] 医药与GLP-1：口服药渗透率加速，多肽制造绿色化成瓶颈 - 减肥药热潮的真正瓶颈不在需求，而在多肽制造的绿色化能力
  \item [R022] 保险与支付：渠道重塑与信任定价 - 市场真正低估的不是卡组织费率，而是“信任即服务”的定价权
  \item [R023] 欧洲汽车与日本IT：结构性变化下的赢家与输家 - 日本汽车业IT投资正在踩刹车，但这恰恰暴露了另一条增长主线
\end{itemize}
\section{逐篇报告摘录}
\subsection{[R001] 美元真正的分裂不在指数里，而在“谁在涨、谁在跌”的背离中}
{\small\color{refgray}归类：美元与外汇：美元走势分裂，日元回流叙事被高估}

如果你只看DXY指数，今年美元似乎涨了1.5\%。但如果把目光放到更全面的贸易加权指数上，美元实际是跌了0.1\%。这不是统计口径的差异，而是全球外汇市场正在发生一场结构性的重新定价——美元不再是一个统一的“强美元”或“弱美元”故事，而是被地缘冲突、能源贸易流和各国央行政策分化撕成了两半。

某外资投行最新发布的全球外汇策略报告，用“Divided Dollar”这个判断贯穿全文。我们认为，这份报告最有价值的洞察不在于它预测了哪个货币对会涨，而在于它提供了一个框架：为什么在同一个宏观环境下，有些货币在涨，有些在跌，以及当环境变化时，资金会流向哪里。

报告用了一个非常精准的对比：DXY指数（欧元权重超过57\%）年初至今上涨约1.5\%，而更全面的美元贸易加权指数却下跌约0.1\%。这意味着，美元的强弱完全取决于你拿它去和谁比。

这背后的驱动力是什么？报告指出，一方面，美国国内的周期性数据仍然强劲——非农就业数据超预期、ISM调查显示增长前景改善与通胀压力并存——这使得利率差持续偏向美元，尤其是相对于欧洲。但另一方面，中国人民币的持续升值（报告认为这一趋势虽然渐进但将超出市场预期）和日元在干预威胁下的“粘性”，使得美元在亚洲和新兴市场方向上反而承压。

更深层的结构是“战争印记”。中东冲突持续近100天，霍尔木兹海峡的能源流动仍然严重受限，这导致贸易条件（terms of trade）在不同经济体之间产生了巨大的分化。能源出口国的货币受益于价格上涨，能源进口国的货币则承受压力。美元在这种分化中扮演了“两面角色”：作为避险货币，它在风险事件中走强；但作为能源进口国的对手方，它又因为能源价格传导而间接承压。

这意味着，只要能源流动没有实质性恢复，美元的分裂状态就不会结束。投资者需要放弃“看多或看空美元”的简单框架，转而关注“美元相对于哪些货币会走强”的精细化判断。

2. 日元的最大风险不是贬值本身，而是市场对“鹰派意外”的定价已过于充分

报告提到，日本央行行长植田和男在6月3日的讲话中似乎暗示了6月加息的可能，这促使该行经济学家将加息预期从7月提前至6月。但报告同时指出一个关键矛盾：市场对鹰派意外的门槛已经非常高——如果日本央行只是传递“渐进加息”的信号，恐怕不足以推动日元持续升值。

为什么？因为日元面临的压力来自两个更根本的宏观因素：一是美国利率“更高更久”的格局仍在持续，且风险仍偏向上行；二是日本国内的财政风险将在未来几个月随着支出计划的明确而重新成为焦点。

更重要的是，报告对“日元回流推升汇率”的预期泼了冷水。除非有政策指令或更有利的利差环境，否则期待企业将海外利润大规模汇回日本来支撑日元，大概率会落空。报告的核心结论是：尽管存在干预风险，他们仍然倾向于将日元作为融资货币来使用。

这里有一个尚未完全回答的问题：如果日本央行真的在6月加息，而市场将其解读为“被迫追赶曲线”而非“主动收紧”，那么加息对日元的支撑能持续多久？报告没有给出明确答案，但这正是需要后续跟踪的关键变量。

3. 一旦冲突缓解，资金将从能源出口国转向能源进口国——但反弹幅度可能低于预期

报告提供了一个非常清晰的“冲突解决交易手册”（Resolution Playbook）。核心逻辑是：过去几个月，市场一直在做多高Beta的能源出口国货币（如澳元、巴西雷亚尔）。但如果冲突缓解，能源价格回落，贸易条件效应的逆转将使得资金从能源出口国转向能源进口国。

报告通过一个二维框架（横轴：对风险情绪变化的敏感度；纵轴：对贸易条件变化的敏感度）来识别受益者。在G10中，瑞典克朗、新西兰元和英镑被视为最佳的“和平交易”标的；在新兴市场中，南非兰特、韩元、波兰兹罗提、智利比索和匈牙利福林被认为最有可能受益。

但报告也给出了一个重要的提醒：即使能源流动完全

[... middle omitted ...]

指出，自美伊冲突以来，卢比的套息水平已经上升，目前不仅高于其他亚洲高收益货币（印尼盾、菲律宾比索），甚至高于南非兰特和墨西哥比索这些传统套息标的。

更重要的是，从该行的估值模型（GSDEER和GSFEER）来看，卢比目前在高息新兴市场货币中属于被低估的。同时，印度央行在最近的政策会议上虽然按兵不动，但推出了一系列措施来缓解资本外流压力——包括对外资投资政府债券的资本利得税和利息税豁免、延长外资可投资债券的期限、以及允许银行发行外币债券和吸收外币存款。

报告的观点是：这些措施将限制卢比的贬值压力，但不会带来明显的升值。他们预计卢比将在当前水平附近“筑顶”，而非“见顶”。这意味着，对于寻求在新兴市场中配置低波动套息仓位的投资者而言，卢比正在成为一个值得重新评估的选项。

但这里有一个尚未验证的假设：如果油价因为冲突缓解而大幅下跌，卢比的套息吸引力会如何变化？报告提到“在出现更明确的冲突解决信号时会考虑将卢比加回多元化套息组合”，但并未量化油价下跌对卢比套息空间的冲击。

\subsection{[R002] 美国利率的真正拐点，取决于一个被忽视的变量}
{\small\color{refgray}归类：宏观与利率：劳动力韧性重塑降息叙事，利率“新地板”已抬升}

这份研报的核心判断值得每一个关注全球资产定价的人仔细琢磨：市场目前对美联储路径的定价，可能低估了一个结构性因素——利率的“新地板”已经抬升，而非短期波动的噪音。

五月非农就业报告的强劲表现，让美国国债收益率在近期找到了一个更高的底部。但这不仅仅是数据本身的问题。真正重要的不是“美联储会不会再加息一次”，而是整个利率分布的形态已经发生了变化。研报明确指出，尽管最终仍预期年底收益率低于当前水平，但近期内，前端的加息定价压力难以消散。

这意味着，投资者需要调整的，不是对某个具体利率水平的预测，而是对整个风险收益框架的重新认知。过去一年里，市场在“衰退交易”和“再通胀交易”之间剧烈摇摆，但这一轮周期可能正在进入一个新的阶段：增长韧性叠加通胀粘性，正在把利率的“地板”钉在一个更高的位置。

研报中一个关键的洞察是，市场对美联储政策路径的定价，目前呈现出一个显著的不对称结构：加息的风险被前置，而降息的空间被压缩。五月就业报告发布后，前端利率的重新定价并非均匀分布，而是朝着一个方向倾斜。

这种不对称性，并非来自市场对美联储单次行动的预期，而是来自对“政策反应函数”本身的重新评估。当经济增长数据持续超预期，而通胀回落速度不及预期时，市场会自然而然地认为，美联储维持高利率的时间会更长，甚至不排除在某个时点再次加息的可能性。研报中提到的“长期暂停”而非“加息周期”的判断，恰恰是这种不对称定价的核心体现。

对于投资者而言，这意味着不能简单地用“加息见顶”的思维来布局。更务实的做法是，承认前端利率在短期内可能维持在一个相对较高的水平，并且任何向下的修正都需要来自数据层面的明确验证。研报中提到的“需要一定数量的数据验证”这一表述，恰恰点明了当前市场定价的脆弱性所在——它依赖于后续数据的配合。

在整份研报中，对1y1y远期利率的分析，可能是最具交易价值的部分。这个期限点自本轮冲突以来，已经成为美国远期曲线上波动最大的环节，累计重新定价幅度高达约110个基点，并且总是在涨跌中扮演领跑角色。

这个现象背后，反映的是市场对政策路径短期不确定性的集中定价。1y1y利率本质上代表了市场对未来一年内一年期利率的预期，它天然地对短期政策变动最为敏感。当市场对美联储下一步行动存在分歧时，这个点就会成为多空博弈的主战场。

研报提出了两个可能的修正路径：要么是通胀信号趋于温和，导致前端利率曲线（1y/1y1y）变平；要么是增长和就业市场更具韧性，使得加息风险的时间点进一步后移，从而导致1y1y/2y1y的曲线变陡。这两种路径对资产价格的含义截然不同。

对于固定收益投资者来说，这个信号意味着，曲线交易可能比方向性交易更具吸引力。在1y1y利率已经大幅波动之后，市场可能正在接近一个关键节点：是继续朝着更陡峭的方向发展，还是开始修正。研报倾向于认为，一段稳定期本身就可以为这个点带来压力缓解，但前提是市场看到足够多的数据来支撑这种稳定。

研报对英国国债的调整值得单独拿出来讲。在将美国10年期国债收益率预测上调至4.4\%的同时，报告也将英国10年期国债的年末预测上调了10个基点至4.5\%。这个调整的直接原因是英国国债近年来表现出的“更高贝塔”。

“更高贝塔”意味着，当全球利率上升时，英国国债的收益率上行幅度更大；而当全球利率下降时，其下行幅度也可能更大。这种特征并非来自英国特有的财政风险或政治风险，而是来自其更高的通胀敏感度和更长的久期结构。

研报中一个有趣的发现是，尽管市场对英国国债的叙事集中在政治和供给风险上，但近期英国国债在资产互换上的表现相对稳定。这暗示，当前英国国债的风险溢价更多是宏观驱动而非供给驱动。这意味着，如果全球利率环境出现变化，英国国债的波动性可能会比美国国债更大。

对于全球资产配置者来说，这个信号提醒我们，不能孤立地看待任

[... middle omitted ...]

场的预期。相反，如果欧洲央行选择在6月加息后“按兵不动”，那么市场可能会迅速将注意力转向更远期的路径。

这个判断的关键在于，欧洲央行面临的是与美国截然不同的经济组合：更弱的经济增长叠加更粘性的通胀。这意味着，欧洲央行的政策选择空间可能比美联储更小。如果欧洲央行被迫在经济增长放缓的背景下继续加息，那么其对欧洲资产价格的冲击可能会比美国更为显著。

研报建议，相比于直接做多欧洲短期利率，做多终端利率（terminal rate）可能更具吸引力，因为经济增长信号正在减弱，而加息本身有助于锚定远期利率。这个策略的本质是，押注欧洲央行加息周期对长期利率的压制作用，而非短期利率的上行空间。

5. 报告尚未完全回答的关键问题：日本央行的“反应函数”会改变吗？

在整份研报中，日本央行的部分可能是最耐人寻味的。尽管日本央行行长植田和男释放了更加鹰派的信号，市场对6月加息的概率定价也大幅上升，但研报发现，曲线在最初变平后又重新变陡，这与典型的鹰派信号导致的“熊市变平”模式并不一致。

\subsection{[R003] 日元回流叙事被高估了：真正驱动汇率的不是日本邮政银行，而是GPIF的未对冲仓位}
{\small\color{refgray}归类：美元与外汇：美元走势分裂，日元回流叙事被高估}

日元回流叙事被高估了：真正驱动汇率的不是日本邮政银行，而是GPIF的未对冲仓位

过去几周，日本邮政银行CEO一句“可能将日本国债持有量翻倍”的表态，迅速点燃了市场对日元大规模回流叙事的想象。10年期日债收益率处于数十年高位，日元汇率仍处于历史低位——这两个事实叠加在一起，看起来确实像是一个“资金回家”的完美剧本。

但某外资投行最新发布的研报给出了一个更冷静的判断：市场对日元回流规模与速度的预期，很可能被过度简化了。真正可能对日元汇率产生实质性影响的，并不是日本邮政银行这类对冲型投资者，而是以GPIF为代表的、持有近6万亿美元海外资产且大部分未对冲的日本机构投资者。而恰恰是后者，当前缺乏大规模回流的充分动机。

这份报告的核心价值不在于否定日元回流的可能性，而在于拆解了“谁在推动回流”以及“什么条件下回流才会真正发生”。它提供了一个区分信号与噪音的分析框架——这对于任何关注日元资产定价、全球利率传导或日本资本流动的决策者而言，都是必要的前提工作。

1. 日本邮政银行的表态更多是利率市场信号，而非汇率市场信号

日本邮政银行持有约5500亿美元的海外资产，主要是债券，是日本除财务省外最大的单一外债持有者。当它的CEO公开表示可能增持日债时，市场本能地将其解读为日元利好。

但这份研报指出一个关键区别：日本邮政银行是一个高度对冲的投资者。它的投资决策主要取决于“对冲后的海外收益率”与“日债收益率”之间的比较，而非简单的名义收益率高低。数据显示，尽管10年期日债收益率已升至3\%以上，但同期全球主要国债收益率也在同步上升，导致日债的相对吸引力并未显著改善。

这意味着，即便日本邮政银行真的增加日债持仓，它对日元汇率的直接影响也相当有限——因为对冲操作已经锁定了汇率风险。真正需要关注的，是它是否在减少海外资产持有、并将资金真正汇回日本。如果是单纯调整日元资产组合，那对利率市场的影响远大于对汇率市场的影响。

报告提供了一个实用的跟踪框架：如果日本邮政银行真的在推动实质性回流，最早会在“国际证券交易”数据中的“银行”分项中显现为海外资产购买减少或卖出；随后，在更滞后的“资金循环统计”数据中，如果“面向中小企业的金融机构”同步出现海外证券卖出和日债买入，才能更确信地将这一流动归因于日本邮政银行。目前，这两个信号都尚未得到确认。

2. GPIF才是真正能撼动日元的关键玩家，但它当前缺乏回流动机

如果日本邮政银行对汇率的影响被高估，那么谁才是真正的变量？答案是GPIF——日本政府养老金投资基金。

GPIF持有近1万亿美元的海外资产，且大部分未对冲。更重要的是，日本的私人养老金在资产配置上往往会跟随GPIF的决策，这意味着GPIF的调整可能引发连锁反应。它的策略评估每五年进行一次，上一次发生在2025年，这意味着重大战略调整在2030年之前不太可能发生。

但GPIF可以在目标配置区间内进行灵活调整。报告测算了一个理论情景：如果GPIF将海外债券持有比例降至目标区间的下限（20\%），并将释放的资金转入日债（使日债比例升至30\%），总规模可达约870亿美元。这个数字听起来可观，但需要指出的是，这样的调整不可能一蹴而就，而且当前的经济条件并不支持这一方向。

关键问题在于GPIF的回报目标。GPIF追求的是名义工资增长率加1.9\%的回报，这意味着它作为未对冲投资者，最关心的是不同资产之间的利率差。尽管过去一年日债收益率上升幅度超过美债，但利差仍然为负。在利率差没有显著改善的情况下，GPIF大规模转向国内债券的动机并不强。

报告明确表示：对于任何协调性的、未对冲的日债回流，现在的吸引力甚至低于今年早些时候——当时美联储降息的前景还没有像现在这样充满不确定性。

当前市场讨论日元回流时，往往聚焦于“日债收益率已经很高了”这个静

[... middle omitted ...]

“现在比年初更不具吸引力”。年初时，市场对美联储降息的预期较强，这意味着美债收益率可能下行、美元可能走弱，从而增强了日元回流的逻辑。但随着通胀数据反复、美联储降息预期推迟，这一逻辑的基础正在被削弱。

对于决策者而言，这意味着判断日元回流时机时，不能只看日债收益率本身，而需要同时跟踪美联储的政策路径、全球利率曲线的变动，以及这些因素如何影响未对冲投资者的“有效利差”。

4. 报告尚未完全回答的关键问题：政策干预的可能性与二阶效应

这份报告提供了一个扎实的分析框架，但有两个关键问题并未完全展开——而这恰恰可能是未来市场波动的来源。

第一个问题是：日本政府或财务省是否会主动推动GPIF调整配置？报告提到GPIF的调整可以“在行政指导下”发生。如果日本政府认为日元过度贬值已经损害了经济，它是否可能通过行政手段引导GPIF减少海外资产、增持国内资产？这种政策驱动的回流与市场驱动的回流，在速度和规模上可能有本质区别。报告对此着墨不多，但这可能是一个需要持续跟踪的政策变量。

\subsection{[R004] 中国经济的真实温度：需求在修复，但供给侧的再定价才是关键}
{\small\color{refgray}归类：中国经济：供给冲击而非需求疲软，消费回暖但地产仍弱}

这份某外资投行在6月初发布的中国高频经济活动追踪报告，表面上是在更新消费、生产、外贸和政策的周度数据。但如果只把它当作一系列图表来读，就错过了这份报告背后真正值得推敲的判断。

报告的结论并不复杂：中国经济正在经历一个“非对称修复”——消费和出行层面的高频指标在改善，但工业生产和投资端的信号却呈现出明显的分化。市场目前讨论最多的是“需求是否见底”，但这份报告揭示了一个更深层的信号：真正需要被重新定价的，不是需求本身，而是供给侧的结构性变化。

为什么这个判断重要？因为过去几个季度，市场对中国经济的预期始终在“弱复苏”和“二次探底”之间摇摆。高频数据虽然短期波动大，但它们能比月度宏观数据更早地捕捉到边际变化。而这份报告恰好提供了一组关键信号：消费信心在回升，但工业品价格和就业指标仍在低位徘徊。这意味着，当前的经济修复并不是全面的，而是有选择性的。对产业决策者和投资者而言，理解这种“选择性”在哪里，比猜测下一个月的GDP数字更有意义。

以下是我们从这份报告中提炼出的四个核心洞察，以及一个尚未被充分讨论的潜在风险。

报告中最亮眼的数据来自消费端。Morning Consult消费者信心指数在最新读数中升至多年新高，这并非一个孤立的信号。与此同时，30城新房成交面积虽然周环比持平，但仍高于去年同期；二手房成交面积也维持在高位。更值得注意的是，国内汽油和柴油价格在6月初被大幅下调，这直接降低了居民的出行成本，对消费意愿形成支撑。

这些数据合在一起意味着什么？市场此前对消费复苏的持续性存疑，但高频率的出行和购房数据表明，居民的实际消费行为正在边际改善，且这种改善并非由短期的价格促销驱动，而是由信心修复和成本下降共同推动。

但这里有一个尚未被充分讨论的结构性问题：消费改善的广度。报告中的交通拥堵指数和国内航班量均出现周度下滑，且航班量仍低于去年同期。这说明消费修复主要集中在“刚性需求”领域（如通勤、购房），而“弹性需求”领域（如旅游、休闲出行）的恢复仍然偏弱。对于消费品和零售行业的决策者而言，这意味着：短期可以押注刚需品类的复苏，但不宜过早对可选消费的全面反弹下注。

生产端的数据呈现出一种典型的“量稳价弱”格局。钢铁需求和生产量在周度层面均小幅回落，但沿海省份的煤炭日耗量却进一步上升，且高于去年同期。港口集装箱吞吐量和船舶货运量也均高于去年同期。

这组数据的矛盾点在于：如果生产活动在放缓，为什么能源消耗在增加？一个可能的解释是，生产结构正在从“高附加值、低能耗”向“低附加值、高能耗”偏移。简单来说，企业为了维持收入规模，不得不增加产量，但产品的议价能力在下降。报告中的化工品价格数据也印证了这一点：硫酸价格周度小幅上涨，但甲醇、尿素等基础化学品价格仅企稳，并未出现趋势性上涨。

这对投资者的启示是：当前工业企业的利润修复，更多依赖“量的扩张”而非“价的提升”。这意味着，那些能够将规模优势转化为成本优势和议价权的龙头企业，将比中小型竞争对手更受益。相反，如果企业仅仅依靠扩大产能来维持收入，其盈利质量将面临持续压力。这一点在钢铁和化工行业尤为明显——行业集中度正在成为决定企业盈利能力的关键变量。

报告中最具政策信号意义的数据是地方政府专项债的发行进度。截至6月初，今年已发行的地方政府专项债规模达到1.5万亿元人民币，这一数字在去年同期仅为约0.8万亿元。同时，PSL（抵押补充贷款）余额在5月继续收缩，表明政策性银行仍在回收流动性。

这两组数据放在一起，揭示了一个清晰的财政逻辑：中央政府正在通过加速专项债发行来对冲地方财政的收缩。1.5万亿元的发行量意味着，基建投资的下行风险正在被政策性对冲工具所覆盖。但这里有一个关键问题尚未被回答：这些资金能否在短期内转化为实物工作量？

报告中没有提供施工开工率或水泥出货量的高频

[... middle omitted ...]

。

但这里存在一个显著的信息缺口。报告中的就业指数是“相对值”，而非“绝对值”。它只能告诉我们就业状况是改善还是恶化，无法告诉我们失业率的真实水平。更重要的是，这份报告没有提供关于“青年失业率”或“灵活就业人口占比”的高频替代指标。对于一个正在经历结构性转型的经济体而言，就业市场的“温度”比“方向”更重要。

这意味着什么？当前市场对消费修复的乐观预期，隐含了一个前提假设：就业市场正在企稳。但如果就业市场的改善仅限于制造业和建筑业（这两个行业受专项债和基建投资直接拉动），而服务业的就业状况仍然疲软，那么消费修复的可持续性将面临挑战。对于投资者而言，这是一个需要持续跟踪的风险因子。未来几周，如果报告能够补充更多关于服务业就业和收入预期的数据，将极大提升其分析框架的完整性。

5. 一个被低估的观察框架：关注“价格信号”而非“数量信号”

综合以上分析，我们认为这份报告提供的最有价值的观察框架，不是去追踪生产量或消费量的绝对水平，而是关注“价格信号”的变化。理由有三点：

\subsection{[R005] 资金流向正在告诉你一个被忽视的结构性切换}
{\small\color{refgray}归类：全球资金流向：美国成唯一增长叙事，工业板块受追捧}

截至6月3日当周，全球股票基金净流入230亿美元，扭转了前一周70亿美元的净流出。同期固定收益基金继续获得超过400亿美元的净流入，货币市场基金资产更是飙升了1220亿美元。这些数字本身并不令人意外——市场情绪在宏观事件后反弹是常见剧本。

某外资投行最新发布的全球资金流向周报揭示了一个被多数投资者低估的信号：资金从“防御性集中”转向“结构性再配置”的拐点可能已经出现。这不仅仅是风险偏好的短期修复，而是一场关于产业周期、区域竞争和资产定价逻辑的深层重置。

当周全球股票基金净流入230亿美元，几乎全部由美国市场贡献。除美国外，发达市场普遍呈现净流出。这已经不是第一次出现这种分化。

过去几年，美国市场一直被视为全球资金的“安全港”——流动性好、制度透明、科技龙头抗跌。但当前这轮流入的驱动力正在发生变化。它不再仅仅是对冲不确定性的防御性配置，而是资金开始相信：美国经济可能在主要经济体中率先完成“软着陆”甚至“不着陆”。

报告数据显示，美国股票基金的净流入在经历了4月份的剧烈波动后，4周移动均值已回到正值区间。这种修复速度比市场预期的更快。更重要的是，同期欧洲、日本等发达市场的资金流向并未同步改善。这意味着资金不是在“抄底全球”，而是在“加注美国”。

对于资产管理者而言，这一信号的含义是：美国市场的相对表现优势可能不是暂时的，而是结构性的。如果全球其他主要经济体无法在短期内提供同样清晰的增长路径，资金向美国集中的趋势还会持续。

2. 工业板块正在成为新的资金聚集地，科技板块的“吸金王”地位被动摇

截至5月底，工业板块的累计资金流入占资产管理规模的比例已达到24\%，在所有行业中遥遥领先。基础设施板块以17\%紧随其后，能源和原材料板块也分别达到15\%和11\%。而科技板块的累计流入占比仅为3\%，甚至低于公用事业和电信。

过去三年，科技板块尤其是AI相关领域，几乎吸走了全球权益资金的全部增量。但这份数据表明，资金的偏好正在发生微妙但实质性的转移。工业、基础设施、原材料这些“旧经济”板块重新获得关注，背后的驱动力可能来自几个方向：全球供应链重构带来的资本开支周期、各国财政刺激向实体基建的倾斜、以及能源转型对传统工业的拉动。

值得警惕的是，科技板块的累计流入占比仅为3\%，而同期科技板块股价表现并不差。这意味着科技股的上涨更多依赖盈利增长和估值扩张，而非资金流入驱动。一旦盈利预期出现边际变化，科技板块的脆弱性可能比市场想象的要高。

对于投资者而言，这一行业切换意味着：过去几年“买科技就行”的策略可能正在失效。工业、基础设施等板块的配置价值正在上升，而科技板块需要更严格的盈利验证。

3. 亚洲资金流向出现严重分化，韩国和台湾在“吸金”，中国和印度在“失血”

截至5月底，韩国股票基金的累计净流入占资产管理规模的比例高达37\%，在所有单一市场中排名第一。台湾以13\%紧随其后。而中国大陆和印度则分别录得23\%和6\%的净流出。

第一，韩国和台湾的半导体产业链正在享受全球AI资本开支的溢出效应。资金流入的持续性与这两个市场的出口数据和产业景气度高度相关。这不是短期交易性资金，而是基于产业周期的战略性配置。

第二，中国大陆和印度同时面临资金流出，但原因截然不同。中国大陆的流出更多反映地缘政治不确定性和国内经济复苏节奏的反复；印度的流出则更多来自估值过高后的获利了结——印度市场在过去两年积累了巨大的涨幅，资金回撤在情理之中。

但报告中的另一张图表提供了更微妙的线索：中国大陆股票基金的资金流出主要来自海外投资者，而来自中国大陆本土的资金仍在净流入。这意味着海外投资者对中国资产的悲观程度可能已经接近极端水平，而本土资金的态度相对乐观。这种背离往往出现在市场底部区域附近，但反转时点仍不确定。

对于全球配置者而言，亚洲不再是“

[... middle omitted ...]

金则出现净流出。这一组合清晰地表明：市场正在为“利率在更长时间内维持高位”做准备。

短期债券的吸引力在于，投资者可以在不承担过多久期风险的情况下获得较高的票息收入。通胀保护债券的需求则反映了市场对通胀黏性的担忧——尽管整体通胀数据在回落，但服务通胀和工资增长仍具有韧性，市场不愿意完全押注通胀会顺利回到目标水平。

长期债券基金的流出则说明，投资者并不相信利率会在短期内大幅下行。如果市场真的预期经济衰退即将到来，长期债券应该是最大的受益者之一。但现实是资金正在避开长期债券，这与“软着陆”甚至“不着陆”的宏观叙事是一致的。

对于债券投资者而言，当前的策略选择变得清晰：与其押注利率下行的时间点，不如接受一个“高利率持续时间更长”的现实，通过短久期和高票息来获取确定性更强的收益。长端利率的波动性仍然较高，不适合大规模配置。

报告的“本周图表”部分关注了英镑的跨境资金流动。数据显示，英镑在2024年以来持续获得净需求，这在一个面临国内政治和财政风险的经济体中显得不太寻常。

\subsection{[R006] A股情绪降温背后，真正的结构性信号尚未被定价}
{\small\color{refgray}归类：A股与港股：情绪降温但未崩溃，结构性信号尚未被定价}

过去两周，A股市场出现了一个值得关注的变化：某外资投行编制的A股情绪指标MSASI加权值从66\%回落至62\%，创业板与全市场日均成交额分别下降9\%和6\%。如果只看这些数字，很容易得出“市场正在冷却”的结论。但这份研报真正想传达的判断，远比情绪降温本身更值得推敲。

*市场真正低估的不是需求疲弱，而是K型分化格局下供给侧的再定价。** 5月制造业PMI滑至50\%的荣枯线，消费出口订单与高耗能生产走弱，但高端制造依然坚挺。这不是一个简单的“好”或“差”的故事，而是一个结构正在剧烈重组的信号。投资者若仅以总量情绪判断方向，很可能错过本轮周期中最关键的定价机会。

报告的核心结论是：短期波动仍将持续，但6-12个月维度上，A股仍有10-12\%的上行空间。这一判断建立在三个尚未被充分讨论的变量之上——二季度盈利改善的路径、政策加速落地的节奏，以及A股相对于离岸市场的结构性优势。

以下是我们从这份研报中提炼出的五个关键洞察，以及一个报告尚未完全回答的问题。

MSASI加权值从66\%降至62\%，4个百分点的降幅本身并不剧烈。但更值得关注的是分项指标的分化。根据报告披露的12个情绪指标权重，30日RSI（权重15\%）和融资余额（权重15\%）是贡献最大的两个因子。5月28日至6月3日期间，30日RSI仅上升1\%，而融资余额基本未变——这意味着市场情绪的下滑并非来自恐慌性抛售，而是来自“参与意愿的边际递减”。

这种递减在成交额上体现得更为清晰。创业板日均成交额降至7790亿元，全市场降至3.02万亿元，分别较上一周期下降9\%和6\%。但股指期货成交额和融资余额并未同步萎缩。这传递出一个微妙信号：交易型资金的热情在消退，但杠杆资金和机构对冲仓位尚未大规模撤退。

对于投资者而言，这意味着当前市场并非系统性风险释放，而是一个“等待催化剂”的状态。情绪指标的下滑更多反映了对宏观数据疲弱的短期反应，而非对长期前景的悲观重估。报告在方法论部分强调，MSASI基于100日移动最小-最大归一化，目的是捕捉中期趋势而非高频噪声。从这一框架看，当前情绪仍处于2024年以来的中等偏上区间，尚未触及需要警惕的极端水平。

5月PMI数据是理解当前市场最关键的宏观拼图。制造业PMI降至50\%，符合市场预期，但分项数据的K型分化才是真正的核心。消费出口订单与高耗能生产走弱，而高端制造依然坚挺——这组对比揭示了一个正在固化的结构：中国经济的增长动能正在从“量的扩张”转向“质的定价”。

报告的经济学家团队判断，4-5月数据指向二季度GDP面临下行压力，政策制定者可能从6月起加速财政落地，以基建为靶向防止增长显著低于4.5\%。这一判断隐含了一个重要假设：政策不是全面刺激，而是定向托底。这意味着受益的行业将高度集中，而非普涨。

从这个逻辑出发，报告对行业的偏好就有了清晰的支撑。上游材料、工业、能源和半导体是报告明确看好的方向。这些板块的共同特征是：受益于高端制造的投资周期、具备定价权或技术壁垒、且与政策定向支持的领域高度重合。相比之下，对金融和红利板块的配置建议是“适度”——这并非否定其防御价值，而是暗示在K型分化格局下，防御性资产的超额收益可能不如结构性成长资产。

5月28日至6月3日期间，南向资金净流入45亿美元，年初至今累计净流入353亿美元。这一数据本身并不令人意外——南向资金持续流入港股已是2024年以来的常态。但报告在整体策略上仍然明确表示“偏好A股优于离岸市场”，这一判断值得深究。

报告给出的理由有三：A股对上游制造和硬科技有更高的敞口、有IPO催化剂、以及国家队支持。但在这三个理由背后，还有一个更微妙的逻辑：离岸市场（港股+ADR）面临的短期风险更为集中。 报告中明确列出了几个近期的波动来源——可能的1季报miss、7月港股IPO解禁、

[... middle omitted ...]

I和融资余额的权重最高，分别达15\%。但有一个细节值得注意：外资被动基金流入的R平方仅为0.5\%，是所有指标中最低的。 这意味着外资流动对A股的解释力在过去的数据样本中几乎可以忽略不计。

这引出一个关键问题：随着北向资金日频数据的停止披露（2024年8月起），以及外资在A股持仓结构的变化，MSASI的预测能力是否正在衰减？报告本身没有直接回答这个问题，但方法论部分提到，数据起始点为2014年，因为“一些市场影响因素在此之前尚未完全发展”。如果市场结构正在发生根本性变化，基于历史R平方的权重分配可能需要重新校准。

对于读者而言，这意味着MSASI仍然是一个有价值的参考框架，但不宜将其视为精确的预测工具。尤其是在外资参与度变化、监管政策调整的背景下，情绪指标的解读需要结合宏观基本面的独立判断。

报告在宏观部分提到，5月PPI环比降至1\%（4月为1.7\%），同比可能因低基数升至4\%。通胀压力受控对政策空间而言是好事，但对某些行业而言，这意味着需求端的价格信号仍然疲弱。

\subsection{[R007] 市场真正低估的不是需求，而是供给侧的再定价}
{\small\color{refgray}归类：能源与大宗：供给侧再定价驱动运价与锂价，中国油价调控成本可控}

过去两个月，全球集装箱航运业经历了一场被多数人误解的行情反弹。上海出口集装箱运价指数（SCFI）自3月以来飙升73\%，跨太平洋和亚欧航线录得两位数周涨幅，船公司普遍征收紧急燃油附加费和旺季附加费。表面上看，这是地缘冲突和旺季前备货驱动的短期脉冲。但某外资投行最新发布的行业深度报告揭示了一个更结构性的判断：当前市场定价的核心驱动力，并非需求端的意外强劲，而是供给侧已经发生且不可逆的再定价。

这份报告基于对亚洲头部船公司1Q26运营数据的系统梳理，结合美国终端市场指标、运力管理动态和监管环境变化，勾勒出一幅“行业盈利能力正在经历结构性抬升”的图景。报告覆盖的标的包括TS Lines、中远海控、东方海外、长荣海运和商船三井等亚洲领先航运商，其核心逻辑值得每一位关注全球贸易和产业周期的决策者仔细推敲。

为什么说这不是又一个“旺季故事”？因为支撑当前运价水平的三个结构性因素——有效运力持续收紧、船公司成本传导能力增强、以及行业竞争格局的优化——都具有超越短期周期的持续性。而市场对此的定价，可能才刚刚开始。

报告中最具洞察力的判断之一，是对“有效运力”与“名义运力”的区分。从全球集装箱船队总量看，运力似乎并不短缺。但港口拥堵、航线绕行（因中东冲突导致的红海绕行好望角）以及船公司主动的空白航行管理，正在吸收大量可用运力。Clarksons的港口拥堵指数在2026年1-5月创下历史新高，这意味着全球相当比例的集装箱船队实际上处于“等待”或“绕行”状态，而非有效运输货物。

这一结构性紧平衡解释了为什么在名义运力增长、需求并未爆发式扩张的背景下，运价却能持续上行。报告指出，当前的有效供给紧张程度远高于名义数据所反映的水平。对于船公司而言，这意味着即使需求增速放缓，运价的下行风险也相对可控，因为供给侧的约束构成了一个坚实的底部。

对行业竞争的启示是：拥有自有船队比例较高的运营商，在这一环境中获得了结构性优势。长荣海运的自有船队比例约为70\%（远高于前十大船公司约57\%的平均水平），TS Lines更是高达83\%。在当前期租费率处于疫情后最高位的背景下（Clarksons期租费率指数报206点），自有船队比例高的公司天然免疫于租金成本上涨，其盈利能力的可持续性远优于依赖租赁运力的同行。

2. 成本传导能力是区分“好公司”与“普通公司”的关键分水岭

报告用大量篇幅分析了船公司应对成本上涨的能力，这恰恰是投资者最容易忽视的维度。中东冲突导致燃油价格大幅上涨（相关航线燃油价格自2月底以来上涨约75\%），同时高低硫油价差扩大至每吨135美元以上，是去年同期的数倍。在这一背景下，船公司能否将成本压力有效传导给客户，直接决定了其盈利能力的稳定性。

报告揭示了一个被低估的结构性优势：亚洲船公司因配备脱硫塔（scrubber）的船舶比例远高于全球平均水平，在高低硫油价差扩大的环境中获得了额外的成本缓冲。长荣海运95\%的运力配备脱硫塔，TS Lines为94\%，中远海控/东方海外为58\%，阳明海运为73\%，而全球行业平均水平仅为47\%。这意味着在同等运价水平下，这些公司的单位运营成本低于全球同行，盈利能力更具韧性。

更关键的是，船公司正在通过征收紧急燃油附加费、调整合同条款等方式，将成本压力系统性地转嫁给货主。报告提到，马士基在1Q26实现了超出预期的盈利，并维持了全年指引，这得益于其强大的成本传导能力和高船舶利用率。这标志着行业定价权的转移：船公司不再是被动接受运价的市场参与者，而是主动管理运价和成本传导的定价者。

对于全球航运需求最大的担忧，往往来自美国终端市场的走弱。但报告引用的物流经理人指数（LMI）显示，2026年5月该指数录得69.5，虽略低于4月的69.9，但仍是2022年初以来第二快的扩张速度。更值得关注的是分项数据

[... middle omitted ...]

事。

对于关注中国出口的读者而言，这一判断尤为重要。中国商务部的表态——中美双方就解决非关税壁垒达成积极共识——进一步强化了全球贸易前景的改善预期。如果中美贸易关系出现边际缓和，跨太平洋航线将直接受益，而这正是当前运价涨幅最大的区域之一。

报告提到，国际海事组织（IMO）关于净零排放框架的投票在2025年10月的特别会议后被推迟，最近结束的MEPC 84会议仍未达成最终决议，下一次特别会议定于2026年底举行。这意味着在2027年之前，全球航运业将处于一个“监管真空期”——没有新的强制性碳排放法规落地，船公司无需为碳排放支付额外成本或进行大规模资本开支。

这一窗口期对船公司盈利能力的正面影响是双重的。一方面，当前无需承担碳税或购买碳配额的成本，运营成本保持低位。另一方面，船公司可以利用这段时间加速对替代燃料船舶和节能技术（ESTs）的投资——目前62\%的船队运力已配备节能技术，较2019年初的35\%大幅提升——为未来的合规做好准备，同时享受当前低成本运营的红利。

\subsection{[R008] 市场真正低估的不是AI需求，而是它如何重塑全球收入分配}
{\small\color{refgray}归类：AI重塑全球收入分配：高收入经济体才是最大输家}

这份由某外资投行发布的研报，试图回答一个被多数讨论忽略的问题：当生成式AI从“生产力工具”进化为“认知替代者”，谁才是真正的受害者？主流叙事说，知识密集型行业首当其冲，低端岗位最先消失，依赖劳动力套利的新兴经济体（印度、印尼、墨西哥）风险最大。但这份报告的数据给出了相反的判断——高收入经济体才是AI冲击下最大的输家，而市场当前的定价恰恰呈现了这种“认知倒挂”。

为什么AI开始大规模替代认知任务时，发达国家的经济基础反而更脆弱？为什么新兴市场的股指跌幅最大，但长期经济损伤却可能更小？这不是一篇关于哪个行业会失业的常规分析，而是一个关于“全球收入金字塔如何被AI重塑”的推演框架。

主流叙事认为，知识产业将遭受最严重打击，而新兴经济体依赖知识产业的劳动力套利，因而风险最高。但这个判断的第一个漏洞在于：信息产业在就业中的占比，在不同经济体之间差异巨大。

研报数据显示，高收入国家中，信息产业（IT、广播电视、电信、信息通信）直接雇佣的劳动力占总量7\%到10\%。而在中等收入国家，这一比例仅为1\%到4\%。以印度为例，即使将间接依赖数字技术的岗位（如金融、医疗、零售）全部纳入，其数字密集型就业占比也远低于OECD国家均值。

这意味着什么？即便AI以最激进的速度替代知识工作者，新兴经济体的就业基础受到的直接冲击在量级上远小于发达国家。这不是说新兴市场没有风险，而是说风险的“体积”被放错了位置。一个经济体如果只有3\%的人在信息产业工作，AI替代掉其中一半，对整体就业的影响是1.5个百分点；而如果这个比例是10\%，替代掉一半就是5个百分点——后者才是真正的结构性冲击。

所以，第一个结论是：AI对就业的破坏力，不是由“知识产业的重要性”决定的，而是由“知识产业在就业中的占比”决定的。从这个维度看，发达国家才是真正暴露在前的群体。

如果说就业占比决定了“谁会被波及”，那么工资差异决定了“谁会被优先替代”。研报给出了一个直观的对比：一个典型的软件工程师在印度的年薪约为1万美元，在美国约为13万美元，在英国约为8万美元。会计师、法务助理、数据分析师——任何AI正在快速渗透的白领岗位——都呈现同样的模式。

这里的逻辑非常直接：用AI替代一个1万美元的工人，节省的成本有限；但替代一个10万美元的工人，对大多数组织而言是变革性的。如果企业最终采用“人在回路”系统，保留一小部分决策者和风险承担者，那么保留哪些人类在经济上最合理？答案显然是那些工资最低但技能足够高的——即新兴市场中相对高薪的那部分人。

这引出了一个反直觉的判断：AI替代的“经济动机”在发达国家远强于新兴市场。发达国家的高工资既是优势，也是负担。当AI能以订阅费替代年薪10万美元的员工时，企业的财务激励几乎不可抗拒。而新兴市场的低工资，反而在一定程度上构成了对AI替代的“天然缓冲”——因为替代的经济收益不够大。

但这个缓冲同时也是陷阱，因为它意味着新兴市场通过劳动力套利实现收入升级的路径正在被堵死。

如果上述分析成立，那么一个矛盾就出现了：为什么自AI工具大规模发布以来，新兴市场（印度、印尼）的股指跌幅最大，而美国、韩国、台湾的市场却大幅上涨？

研报给出了一个清醒的解释：市场并非完全非理性，但它是短视的。当前市场的定价逻辑是：谁是AI的直接受益者？答案是那些生产半导体、设计芯片、拥有基础模型的经济体——美国、韩国、台湾。印度和印尼既不是AI的供应方，也没有纯正的AI公司，它们是被颠覆的一方。因此，在缺乏受益标的的情况下，这些市场的股指因IT服务板块（占指数约5\%至10\%）的拖累而下跌，是合理的。

但问题在于，市场定价的是未来几个季度，而不是未来几十年。当市场聚焦于“谁在AI供应链中赚钱”时，它忽略了一个更深层的长期问题：AI对最终需求的破坏——即谁的收入会被消灭

[... middle omitted ...]

设：如果让AI自由发展，每个企业都会加速采用，直到AI替代的岗位数量超过它能创造的新岗位。在那个临界点，真正的经济金字塔重构将发生。

研报的估算显示，在AI饱和的世界里，中等收入国家的国民收入平均减少约10\%，而高收入国家的国民收入减少约22\%。这个差距来自两个因素：一是高收入经济体对知识产业的依赖度更高，二是高收入经济体的工资水平更高，使得AI替代的经济收益更大。

但更重要的不是总量，而是分配。报告认为，AI带来的收入损失将不成比例地落在最富有的1\%之后的“次富9\%”群体——即高收入国家中那些年收入在10万到20万美元之间的专业人士。这恰恰是过去二十年全球化中受益最多的群体，也是消费升级的主要驱动力。

这意味着什么？AI可能在两个维度上成为“均衡器”：一是缩小高收入与中等收入国家之间的收入差距（因为发达国家损失更大）；二是扁平化各国国内除最富1\%之外的收入分布（因为中间层受损最重）。最终，全球可能形成一个极其富有的顶层1\%，以及一个收入更加均等化的其余99\%。

\subsection{[R009] 就业市场比通胀更值得关注的信号：美联储的降息叙事正在被重新定价}
{\small\color{refgray}归类：宏观与利率：劳动力韧性重塑降息叙事，利率“新地板”已抬升}

就业市场比通胀更值得关注的信号：美联储的降息叙事正在被重新定价

市场对美联储降息的预期，可能正建立在一个正在被数据推翻的假设上：劳动力市场会显著放缓。但最新发布的5月就业报告表明，实际情况恰恰相反。

非农就业人数增长17.2万，远高于市场预期的8.8万，且前两个月数据被上修。失业率稳定在4.3\%，而衡量失业的多个指标在5月出现回落，逆转了4月的恶化趋势。这不是一个正在冷却的劳动力市场——这是一个在政策利率维持高位的情况下，依然表现出韧性的市场。

某外资投行在最新研报中明确指出：这份报告应让美联储官员继续聚焦通胀风险。其核心判断是，当前政策立场并未对经济构成实质性抑制。这一结论对资产定价的含义是深远的——如果劳动力市场不降温，通胀回落至2\%目标的路径就更难被确认，降息的时间窗口可能比市场预期的更远。

以下是我们从这份研报中提炼的五个核心洞察，以及在数据背后值得追问的关键问题。

非农就业17.2万的增量本身已经超出预期，但真正重要的信号藏在结构里。研报使用的扩散指数显示，就业增长的广度在近期持续改善，周期性行业的表现尤为突出。这意味着增长不是由少数几个板块拉动的，而是具有更广泛的产业基础。

具体来看，休闲酒店业增长了7万，部分受世界杯临时性因素推动；地方政府就业增加了5.5万。这些看似有“噪音”的细节，恰恰说明劳动力市场的韧性并非虚胖。即便剔除这些一次性因素，私营部门就业的3个月移动均值也已攀升至2024年3月以来的最高水平。

对于投资者而言，这意味着一个关键判断：劳动力市场的韧性正在从“结构性短缺”转向“周期性复苏”。如果后者成立，那么工资增长和消费支出的下行空间将比市场预期的更有限。

4月失业率数据曾引发市场对劳动力市场恶化的担忧。但5月数据给出了相反的信号：衡量失业的多个指标——包括初请失业金人数、JOLTS裁员率、短期失业率——都在5月出现回落。研报特别指出，求职成功率（job-finding rate）虽然仍显疲弱，但已较4月有所回升。

这意味着什么？4月的数据可能是一次统计噪音，而非趋势反转。失业率稳定在4.3\%，同时失业者找到工作的速度在改善，这组组合指向的是一个“有摩擦但不出清”的劳动力市场。对于美联储而言，这恰恰是最难处理的状态：既不足以触发降息，也不至于引发衰退恐慌。

3. 核心CPI可能放缓，但核心PCE却在加速——市场在关注错误的通胀指标

研报预测5月核心CPI环比增速将从4月的0.376\%放缓至0.183\%，主要受租金相关技术性调整的暂时性影响消退。但与此同时，核心PCE通胀预计将从4月的0.239\%加速至0.327\%，同比增速进一步升至3.4\%。

这两组数字的分歧源于统计口径差异：PCE中的投资管理服务价格和机票价格（均来自PPI相关分项）在5月表现强劲。而美联储明确以PCE作为通胀目标。

这里有一个容易被忽略的细节：即便核心CPI放缓，也只是因为4月的一次性因素消退。研报明确指出，核心商品通胀仍保持小幅正增长，科技相关价格压力正在抵消关税影响的消退。更值得警惕的是，ISM服务业价格支付指数已升至2022年中以来最高，企业报告的成本上升范围涵盖能源、货运、包装材料、食品、化肥、建筑材料甚至劳动力。

研报特别点名了达拉斯联储主席Logan的发言。她明确表示，当前经济状况表明“货币政策并未抑制经济”，并“越来越担心今年晚些时候可能需要提高利率”。这是近期美联储官员中，第一个无条件主张加息的。

克利夫兰联储主席Hammack也维持鹰派立场，认为5月就业报告确认了劳动力市场大致平衡。她虽然没有直接主张加息，但明确表示维持利率稳定是合理的。

只有纽约联储主席Williams仍属鸽派，认为通胀不会持续，且没有明显理由修改前瞻指引。但他目前是少数派。

一个关键信号是：

[... middle omitted ...]

的传导渠道已经打开。企业正在将上升的投入成本转嫁给下游，而劳动力市场的韧性又为这一传导提供了支撑。即便关税本身没有大幅升级，伊朗战争引发的能源冲击、货运成本上升、金属价格上涨等，都在通过供应链层层叠加。

研报估算，从国家层面的IEEPA关税转向普遍的10\% 122条关税，使有效关税率下降了约1.5个百分点。而301关税的全面实施很可能将逆转大部分降幅，使有效关税率从当前的约7\%回升至8.5\%。这不是一次性的冲击，而是一个持续抬升的过程。

报告尚未完全回答的关键问题：劳动力市场的“韧性”能持续多久？

研报对劳动力市场的判断是“健康但不过热”——失业率稳定，工资增速放缓速度快于预期。但有一个关键假设尚未充分验证：当前就业增长的韧性有多少来自一次性因素（世界杯、地方政府招聘），多少来自结构性改善？

报告本身也承认，休闲酒店业的7万增长可能部分受世界杯推动，而地方政府就业的增长是否可持续，取决于州和地方政府的财政状况。如果这些临时性因素消退，就业增长能否维持当前速度？

\subsection{[R010] 欧洲汽车市场真正被低估的风险，不是中国份额，而是欧洲的利润优先策略}
{\small\color{refgray}归类：欧洲汽车与日本IT：结构性变化下的赢家与输家}

欧洲汽车市场真正被低估的风险，不是中国份额，而是欧洲的利润优先策略

欧洲正在经历一场它不愿承认的结构性变化。中国汽车制造商在欧洲的市场份额在过去12个月内翻了一倍，达到约6.8\%。这个数字本身并不惊人，真正值得关注的，是这背后两股力量的分化：中国车企在加速渗透，而欧洲车企却在主动选择“不反击”。

这不是一个关于“谁将赢得市场”的故事，而是一个关于“谁愿意为市场份额牺牲利润”的博弈。某外资投行最新发布的全球汽车研报给出了一个反直觉的判断：市场过度关注了中国的进攻，却低估了欧洲的防守策略——后者可能才是决定未来五年行业利润分配的核心变量。

这份报告的核心洞察在于：中国车企的崛起并非靠简单的价格战，而是产品力的系统性提升；欧洲车企并非无力反击，而是选择了保护利润而非份额。这两种策略的碰撞，正在重塑欧洲汽车业的竞争逻辑。

中国车企在欧洲的市场份额从2023年初的约1\%攀升至2026年4月的近7\%，这一增速远超历史上任何新进入者。以英国市场为例，奇瑞集团和比亚迪在极短时间内就达到了3-4\%的市场份额，而丰田和现代当年达到这一水平耗费了数倍的时间。

报告特别强调了一个关键判断：中国车企的竞争力并非来自价格战。在德国市场，中国品牌的定价基本处于各细分市场的中位区间，并未出现系统性低价倾销。真正驱动消费者选择的，是产品本身的“物超所值”——更丰富的配置、更精致的工艺、更先进的智能化功能。

这一点在报告中得到了实地验证。该投行的分析师团队在2025年北京车展和英国经销商走访中发现，中国汽车的产品质量已经不再是“够用”，而是“更好”。这背后是中国国内市场的“完美竞争”逻辑：大量生产资质和地方政府资金支持，迫使车企在配置和品质上不断内卷，最终催生出全球范围内最具竞争力的产品矩阵。

反补贴关税并未有效阻挡这一趋势。欧盟对中国产电动汽车加征的关税，中国车企似乎选择了直接吸收，而不是转嫁给消费者。这意味着，关税壁垒的有效性正在被中国车企的成本优势和利润空间所削弱。

面对中国车企的攻势，欧洲车企的反应出乎许多人意料——它们没有选择在价格上硬碰硬，而是明确将保护利润和现金流作为首要目标。

报告引用了多家欧洲车企高管的公开表态。雷诺CEO明确表示“中国车企应该与欧洲合作、在欧洲投资生产”，而非直接对抗。Stellantis则强调其品牌组合优势是对抗中国新进入者的核心武器。福特CEO更是直言“我们知道自己正处在一场生死之战中”。

但这些表态背后，是一个更现实的商业逻辑：欧洲汽车市场本身就是一个利润微薄的市场。产能过剩、僵化的劳动法、小型车主导的产品结构，使得在这个市场赚钱本就困难。在这样的背景下，与中国车企展开全面的价格战，只会加速利润率的崩塌。

因此，欧洲车企的策略是双管齐下：一方面通过采用磷酸铁锂（LFP）电池等低成本技术降低可变成本，另一方面利用老龄化劳动力自然减员来压缩固定成本。这是一种“防守型”策略——不是不反击，而是选择在自己有优势的领域（品牌、渠道、服务）上竞争，而不是在价格上硬拼。

报告的数据揭示了另一个重要洞察：中国车企的市场份额增长并非来自欧洲本土车企，而是主要蚕食了其他亚洲品牌和部分美系品牌的市场。

在整体市场份额变化中，Stellantis是最大的输家，过去三年流失了约3.9个百分点的份额。其次是福特（-1.5\%）、起亚（-0.8\%）和现代（-0.5\%）。而大众集团反而实现了1.5个百分点的增长，雷诺和沃尔沃也略有上升。

这一格局意味着：中国车企的冲击波并非均匀分布。对于部分欧洲本土车企而言，中国品牌的崛起可能反而“帮”它们挤走了日韩品牌和福特等竞争对手。这或许解释了为何一些欧洲车企并未表现出强烈的反击意愿——只要自己的份额没有受到直接冲击，它们更愿意坐视竞争对手的退场。

在纯电动车领域，竞争

[... middle omitted ...]

公司车队采购，残值管理至关重要。但如果中国车企的本地化生产（如比亚迪在欧洲建厂）进一步降低成本，或者它们开始更激进地定价，欧洲车企的利润优先策略将面临严峻考验。

另一个未解之谜是：欧洲车企的“成本削减”能走多远？通过采用LFP电池降低可变成本，通过自然减员降低固定成本，这些措施都有其极限。当成本削减空间耗尽，而中国车企仍在不懈推进时，欧洲车企是否会被迫转向价格战？

此外，报告提到了一个值得警惕的趋势：中国汽车正在向更大尺寸演化，这可能使其与欧洲市场的需求错位。但这一判断仍需验证——如果中国车企能够同时提供多种尺寸的产品，这一“错位”就可能不复存在。

基于报告的分析逻辑，我们可以提炼出三个核心观察框架，用于跟踪这一趋势的演进：

*第一，关注中国车企的定价策略是否出现系统性变化。** 目前中国车企在德国的定价基本处于各细分市场的中位区间。如果它们开始主动降价，意味着竞争进入新阶段。反之，如果它们继续维持当前定价，说明产品力本身足以支撑增长，这对欧洲车企的威胁更大。

\subsection{[R011] 中国油价调控的真正成本，比市场担心的低得多}
{\small\color{refgray}归类：能源与大宗：供给侧再定价驱动运价与锂价，中国油价调控成本可控}

当霍尔木兹海峡的阴影笼罩全球能源市场，3月以来布伦特原油在每桶80-130美元区间剧烈波动时，中国决策者面临一个老问题：如何在外部冲击与内部稳定之间找到平衡。市场普遍关注油价上涨对通胀和消费的冲击，但这份来自某外资投行的研报给出了一个反直觉的核心判断：中国零售燃油价格管理的财政成本是可控的，年化不到GDP的0.3\%，且这一成本正在被结构性变量——尤其是新能源汽车的加速渗透——所稀释。

真正值得关注的不是短期油价波动本身，而是中国在应对这场危机中暴露出的定价机制弹性、需求侧的结构性转移，以及政策空间的实际边界。这份报告的深度在于，它没有停留在描述“政府干预了什么”，而是量化了“干预的代价有多大”，并指出了“哪些成本被市场高估了”。

1. 定价机制的核心不是“补贴”，而是一个精密的非线性传导器

中国成品油定价机制常被外界简单理解为“政府补贴油价”，但报告揭示的框架远比这复杂。其核心是一个以国际原油价格每桶40-130美元为边界的非线性传导系统。

当布伦特原油低于40美元时，国内零售价格不再下调，价差部分进入“风险准备金”，用于节能减碳和油品升级。当高于130美元时，零售价格冻结，政府通过转移支付补偿炼油商损失。而在40-130美元的正常区间内，每桶80美元又构成一个关键阈值：低于80美元时，国际价格变动几乎全额传导；高于80美元时，政府启动“部分传导”机制，要求炼油商通过压缩加工利润来吸收部分冲击。

这个机制的精妙之处在于，它天然具备“自动稳定器”功能。在油价温和上涨时，市场化传导维持了价格信号的有效性；在油价极端波动时，行政干预介入保护实体经济。报告特别指出，自中东冲突以来，国家发改委在七次调价中，实际累计上调幅度比模型隐含的理论值低了约47\%，相当于只传导了国际涨幅的一半。

这意味着什么？对于下游企业和消费者，这是一个“成本缓释器”，但代价由上游炼油环节承担。对于投资者，这意味着中石化、中石油等国有炼油商的利润表，实质上承载了部分宏观稳定职能——这是评估这些公司估值时不能忽略的制度性折价。

报告中最具洞察力的发现之一，是需求对价格的反应弹性正在发生结构性变化。历史数据显示，国内零售油价每上涨10\%，成品油零售量大约下降5\%。但2026年3-4月的实际数据却呈现出一个显著偏离：零售均价同比上涨17\%的同时，零售量同比下降了17\%。

这意味着需求弹性从约0.5跃升至约1.0，翻了一倍。报告给出的解释是，非石油能源供应和新能源汽车的快速普及，提供了比以往油价上涨周期更广泛的替代选择。这不是一个短期现象，而是一个结构性拐点。

数据支撑了这一判断：新能源汽车渗透率从2016年的1\%上升到2025年的54\%，2026年5月进一步攀升至63\%。五一假期期间，高速公路电动车充电需求同比增长52.8\%，而私家车出行仅增长2.6\%。这些数字共同指向一个结论：中国经济的石油需求强度正在加速下降。

对于关注能源转型的投资者，这意味着传统石油需求的“价格弹性”正在被“技术替代弹性”所覆盖。油价越高，替代速度越快，形成一个自我强化的负反馈循环。这反过来又降低了长期油价冲击对中国经济的潜在伤害——这是一个市场尚未充分定价的宏观韧性来源。

市场对油价干预的担忧，往往集中在“财政补贴会拖累政府预算”。报告用数据打消了这个疑虑：根据估算，政府干预的年化成本仅占GDP的0.3\%左右，并且这一成本还可以通过库存估值收益和上游盈利能力增强来进一步对冲。

但“财政成本可控”不等于“没有成本”。成本只是被转移到了不同的承担者身上。报告明确指出，当国际油价高于80美元/桶时，国有炼油商被要求吸收部分涨价压力，这意味着它们的加工利润被压缩。这是一种隐性的“准财政支出”，不体现在政府预算中，但体现在国有企业的利润表中。

此外，报告还揭示

[... middle omitted ...]

调整正在重塑贸易流向，但存在一个尚未回答的关键问题

除了价格干预，中国还采取了供应侧的应对措施：临时收紧成品油出口、从俄罗斯和拉丁美洲大幅增加原油进口。报告估算，3-4月中国原油进口总量同比下降11\%，但结构发生了剧烈变化：中东进口暴跌38\%，而俄罗斯进口增长13\%，拉丁美洲进口增长64\%。

这种“东退西进”的贸易流向调整，在短期内缓解了供应缺口，但报告没有完全展开的一个关键问题是：这种替代的可持续性如何？中东原油的API度和含硫量与俄罗斯、拉美原油存在差异，这可能影响炼油厂的运营效率和产品产出结构。此外，长期依赖非中东原油，意味着中国在能源安全上需要承担不同的地缘政治成本和运输风险。这些二阶效应，报告没有给出量化评估，但却是产业决策者必须纳入考量的变量。

另一个悬而未决的问题是，政府正在利用这次冲击加速“新三样”——电动汽车、锂电池、太阳能——的政策支持。这无疑会降低经济的长期化石燃料依赖，但加速的节奏是否会在油价回落后面临政策回摆？这是一个值得跟踪的观察点。

\subsection{[R012] 锂价短期回调不是反转，真正被低估的是CATL产能释放带来的需求硬缺口}
{\small\color{refgray}归类：能源与大宗：供给侧再定价驱动运价与锂价，中国油价调控成本可控}

锂价短期回调不是反转，真正被低估的是CATL产能释放带来的需求硬缺口

2026年6月初，碳酸锂价格在触及20万元/吨阻力位后出现回调，现货报价从5月底的17.55万元/吨回落至16.83万元/吨。市场情绪随之转弱，部分投资者开始质疑锂价上行动力是否已经耗尽。

但某外资投行最新发布的研报给出了一个与市场直觉相反的判断：基本面并未松动，锂价从当前水平仍有上行空间。支撑这一判断的核心变量，并非市场热议的供给侧减产或政策刺激，而是一个被大多数人忽视的硬需求增量——CATL的新增产能将在2026年8至9月前投产，届时将带来单月约5000至6000吨的增量锂需求。

这不是一个模糊的“需求复苏”叙事，而是一个可以量化的、时间节点明确的供需错配窗口。市场当前的低迷情绪，恰恰为理解这一结构性变化提供了时间差。

截至6月4日当周，碳酸锂和氢氧化锂的现货均价分别环比下跌4.2\%和5.5\%。从表面看，价格走弱似乎印证了市场对需求疲软的担忧。但拆解周度生产与库存数据后，结论恰恰相反。

产量方面，当周中国碳酸锂产量环比增长3\%至26344吨。其中，盐湖提锂产出环比增长1\%，锂云母提锂环比下降4\%，锂辉石提锂环比增长5\%，回收提锂环比增长3\%。产量增长并未引发累库，反而库存出现了边际下降。

总库存当周为98786吨，环比减少630吨，降幅约1\%。更关键的是库存结构的变化：下游正极材料厂商的库存环比增长7\%至46429吨，冶炼厂库存环比下降2\%至16615吨，而其他环节（主要是电池厂和贸易商）的库存环比大幅下降8\%至35742吨。

这意味着，价格回调期间，下游正在主动补库，而上游和中游的库存压力在减轻。这种库存转移本身就是一个价格支撑信号——如果下游对后市悲观，他们不会在价格下跌时增加库存。当前的行为更接近“逢低买入”而非“恐慌抛售”。

研报中最值得关注的判断，并非对锂价的短期预测，而是对CATL2026年产量的上修：预计CATL全年产量将达到1.2TWh，同比增长61\%。这一数字本身已经超出市场一致预期，但真正重要的不是总量，而是产能投放的节奏。

报告明确指出，CATL的新增产能将在2026年8至9月前陆续投产。这意味着，从三季度开始，CATL将有一个明显的产量跃升。按照当前锂盐单耗测算，这一新增产能将在投产当月带来约5000至6000吨的增量锂需求。

这5000至6000吨单月增量意味着什么？对比当前周度产量约26000吨、月产量约10万吨的规模，这个增量相当于现有月产量的5\%至6\%。在一个供需本就偏紧的市场中，5\%的需求增量足以改变价格走向。

更重要的是，这个需求增量不是假设性的——CATL的产能建设进度是确定的，投产时间表是明确的。市场目前对锂价的悲观情绪，很大程度上是因为没有将这个确定性增量纳入定价框架。

从研报提供的月度库存图表可以看出，自2025年下半年以来，碳酸锂总库存从约10.5万吨逐步攀升至2025年底的约13.5万吨，随后在2026年初回落至约6.5万吨。这一去库幅度值得深思。

2026年初的去库并非需求突然爆发，而是供给端出现了预期外的收缩。2025年下半年，冶炼厂库存一度攀升至约6.5万吨的历史高位，而到2026年初已回落至约2.5万吨。这说明上游在价格低迷期主动减产、去库存，为后续价格上涨奠定了基础。

当前库存水平约9.9万吨，处于历史中位区间。但考虑到CATL即将释放的需求增量，以及下游正在进行的主动补库，这个库存水平并不高。一个简单的测算：如果CATL在8至9月投产，单月增量需求5000至6000吨，即便其他需求保持不变，现有库存也只够支撑约4至5个月的增量消耗。这还不考虑其他电池厂和整车厂的同步扩产。

4. 报告尚未完全回答的关键问题：锂辉石精矿的定价机制是否会被打破

研报中提

[... middle omitted ...]

的假设：锂辉石精矿的供给能否快速响应价差信号。澳洲和非洲的新增矿山投产周期通常需要12至18个月，短期供给弹性有限。如果精矿供给无法跟上，加工利润的扩大可能只是昙花一现，最终将通过锂盐价格上涨来重新平衡。

这个问题的答案，将决定锂价上行的空间和持续性。如果精矿供给弹性充足，锂价可能在25万元/吨附近遇到阻力；如果精矿供给受限，锂价可能突破25万元/吨并向更高区间运行。

基于上述分析，投资者不应将锂价本身作为判断行情的唯一指标。三个先行信号比价格更能反映供需的真实变化：

第一个信号是CATL的月度产量数据。如果CATL在7至8月的实际产量开始向1.2TWh的年化目标靠拢，那么锂需求增量将从预期变为现实，这是最直接的催化剂。

第二个信号是冶炼厂库存的周度变化。如果冶炼厂库存持续下降且下游库存持续上升，说明库存转移正在加速，价格上行压力在积累。

第三个信号是锂辉石精矿的现货报价。如果精矿价格开始跟随锂盐价格上涨，说明上游供给端也在收紧，供需错配正在向全产业链传导。

\subsection{[R013] 市场真正低估的不是需求，而是供给侧的再定价}
{\small\color{refgray}归类：中国经济：供给冲击而非需求疲软，消费回暖但地产仍弱}

中国4月工业增加值同比增速从3月的5.7\%骤降至4.1\%，远低于市场预期的6.0\%，甚至低于某外资投行自己更谨慎的5.2\%预测。这一读数创下近三年来的最低月度纪录。市场普遍将其解读为需求疲软，但这份研报提供了一个截然不同的叙事框架：供给冲击才是主因。

报告的核心判断是，中东能源设施受损与霍尔木兹海峡封锁导致的原材料供应中断，已经对中国工业生产造成了结构性伤害。据估算，4月工业增加值增速骤降中，约有一半来自石油和化工两个子行业的直接与间接拖累。这两个行业合计占工业总产出的15\%，却贡献了原材料行业40\%的产出。

更关键的是，随着美国总统特朗普警告海峡封锁可能持续到9月，这种供给瓶颈在未来几个月内不会自然消退。研报预计5月工业增加值同比增速将维持在4\%左右的低位，并将第二季度GDP增速预测下调至4.1\%，远低于市场共识的4.7\%。

这并非一次短暂扰动，而是一场供给侧压力测试。以下五个层次的分析将揭示：为什么这次供给冲击比表面数据更严重，以及它如何重塑我们对中国经济短期走势的判断框架。

1. 石油化工行业的产出收缩并非需求问题，而是原材料断供的直接结果

4月工业增加值的行业分解揭示了一个清晰的传导链条。石油、煤炭及其他燃料加工业的产出增速从3月的8.1\%骤降至-0.9\%，这是自2024年8月以来的最低水平。原油加工量同比收缩5.8\%，较3月的-2.2\%进一步恶化。沥青产出量同比暴跌40.1\%，而3月这一数字是-19.7\%。

化工原材料及制品业的产出增速也从3月的9.0\%放缓至5.3\%。化学纤维制造业的产出增速降至2.2\%，为2023年3月以来最低。这些数据指向同一个问题：原材料进口断供。

报告提供了一组关键数据：4月中国原油进口量同比暴跌20.0\%，降至2022年8月以来最低水平。来自俄罗斯的替代供应似乎不足以填补缺口——3月和4月从俄罗斯的原油进口增速分别放缓至13.8\%和11.3\%，远低于1-2月40.9\%的激增幅度。

石脑油进口量同比暴跌51.3\%，而中国约17\%的石脑油依赖进口，其中约40\%来自波斯湾。液化天然气进口量同比下降23.1\%，创2018年4月以来最低月度读数，原因是来自卡塔尔的进口量同比暴跌99.4\%——伊朗导弹反击严重破坏了卡塔尔的拉斯拉凡出口设施。

这些数字意味着什么？不是需求不足导致企业不愿意生产，而是企业想生产但没有原料。这种供给驱动的产出收缩，与需求疲软驱动的产出收缩在政策应对上有着根本区别。

2. 硫磺短缺正在通过硫酸传导链威胁化肥和农业，这是报告尚未完全展开的二阶冲击

在众多原材料中，硫磺的供应紧张尤为值得关注。4月中国硫磺进口量同比暴跌72.4\%，降至2008年10月以来最低水平。中国对硫磺的进口依赖度约为50\%，其中约55\%来自中东。

更令人担忧的是，4月底至6月初，中国主要港口的硫磺库存进一步下降26\%，此前3-4月已下降了28\%。硫磺价格自2月底以来已飙升94.4\%。

全球80\%-90\%的硫磺最终用于生产硫酸。而中国是全球最大的硫酸生产国和出口国。硫酸总消费量的55\%-60\%用于化肥生产，主要是磷肥。

这里有一个尚未被充分讨论的二阶效应：如果硫磺短缺持续，硫酸产量下降将直接传导至磷肥生产，进而影响农业生产成本和粮食供应。报告提到了这一点，但没有量化潜在影响。考虑到中国在全球化肥供应链中的核心地位，这一传导链条可能产生远超工业领域的连锁反应。

报告尚未完全回答的问题是：中国的战略储备能否填补这一缺口？储备规模、释放节奏以及替代来源的获取能力，将决定这一冲击的持续时间和严重程度。

研报提供的5月高频数据令人不安。山东地炼（小型独立炼厂，掌握中国约25\%的炼油产能）的开工率从4月底的60.0\%降至5月底的54.7\%，6月初进一步降至53.6\%。沥青工

[... middle omitted ...]

但研报指出，如果调整供应商配送时间分项指数（这是一个反向指标）的扭曲效应，综合PMI可能已经跌破50，更接近其预测的49.9。

更值得关注的是，国家统计局指出，5月PMI的主要拖累因素包括石油、煤炭及其他燃料加工业，化学纤维及橡胶塑料制品业，这些行业的生产和新订单分项指数均低于50荣枯线。这意味着供给和需求两端都在走弱——但根本原因仍然是供给侧的原材料短缺。

这些高频数据传递的信息是：4月并非一次性冲击，供给压力正在持续并向更长的产业链条扩散。石油化工行业占原材料行业产出的40\%，其产出收缩将通过投入产出关系产生乘数效应，放大对整体经济的拖累。

在整体工业增加值疲软的背景下，有一个亮点值得注意：计算机、通信设备及其他电子设备制造业的产出增速从3月的12.5\%加速至4月的15.6\%，为整体工业增加值贡献了0.4个百分点的同比拉动。

但研报揭示了一个被忽视的隐含条件：如果没有科技相关产业的拉动，4月工业增加值增速将跌破4\%。换言之，制造业的真实状况比表面数据更差。

\subsection{[R014] 市场真正低估的不是地缘风险，而是中国药企已不可替代的全球研发节点地位}
{\small\color{refgray}归类：A股与港股：情绪降温但未崩溃，结构性信号尚未被定价}

市场真正低估的不是地缘风险，而是中国药企已不可替代的全球研发节点地位

2026年6月初，港股医药板块在ASCO数据公布后依然呈现“卖事实”的弱势表现。表面看，市场是在消化地缘政治带来的监管不确定性——美国国会接连提出将生物技术纳入COINS法案、引入BINSA法案，试图限制美国资本与中国生物科技公司的交易。但一份来自某外资投行的最新研报提出了一个值得反复推敲的判断：市场可能严重高估了这些法案的实际影响，同时严重低估了中国药企在全球创新药研发链条中已经形成的结构性位置。

这不是一个情绪层面的判断，而是一个基于全球R\&D投入分布、临床资产供给格局、以及交易对手方利益结构的硬逻辑。报告的核心主张是，如果美国真的选择“脱钩”，受损更大的可能不是中国，而是美国自身的生物技术产业。这个结论听起来有些反直觉，但报告用三组数据链条支撑了这一判断。

我们详细拆解这份报告的核心逻辑，并沿着其分析框架，探讨一个更深层次的问题：这轮地缘博弈，真正考验的是什么？

1. 中国已不是“低成本替代”，而是全球最大的创新药早期资产供给方

许多投资者对中国的印象仍停留在“仿制药大国”或“低成本CRO基地”。但报告提供的数据揭示了另一个现实：中国在创新药早期研发阶段的资产供给能力，已经全球领先。

2020年至2025年，中国制药研发支出从260亿美元增长至约390亿美元，年复合增长率8.5\%，占全球制药研发支出的比例稳定在11\%至13\%之间。但更关键的是产出效率——自2020年以来，中国每年进入临床试验的创新药数量已经超过美国，成为全球第一。2024年，中国有722个创新药进入临床试验，而美国为457个。2025年，这一数字进一步拉大到827比482。

这些数字意味着什么？中国已经从一个“跟随者”变成了全球新药候选分子的最大“矿场”。在ADC、多特异性抗体等前沿领域，中国企业的资产积累尤为突出。任何一家全球药企，如果试图绕过中国进行早期管线布局，将面临候选分子数量锐减、同质化竞争加剧的风险。

这不是“成本优势”的问题，而是“供给结构性依赖”的问题。报告指出，美国的生物技术产业如果拒绝与中国合作，将错过大量有潜力的药物候选分子。这不仅是商业判断，更是产业逻辑。

2. 法案对美国的伤害可能远大于对中国——交易对手利益才是真正的阻力

BINSA法案的逻辑是“防止美国资本和技术流向中国生物技术领域”。但报告揭示了一个关键的反向机制：美国生物技术产业同样高度依赖来自中国的资产输入。

2020年至2026年5月，中国生物科技公司的对外授权交易中，美国买家占据了42\%的交易数量和57\%的前期付款金额。这意味着，美国药企是中国创新药资产最大的“进口商”。如果法案落地，这些交易不会消失，而是会流向欧洲、亚洲其他地区的买家。

报告进一步指出，来自美国本土产业界的强烈反对是可以预见的。一方面，限制合作意味着美国患者可能无法及时获得来自中国的创新疗法；另一方面，在公众和政府要求降低药价的背景下，切断与中国的高效研发合作只会推高美国药企的研发成本。这种“自损八百”的政策，很难获得产业界的支持。

更重要的是，立法过程本身将是漫长且充满博弈的。报告认为，市场对消息面的反应可能过度了。2026年前五个月，中国药企对外授权交易的前期付款总额已达48.33亿美元，同比增长93\%，相当于2025年全年的68\%。其中5月单月就贡献了12.5亿美元，主要由恒瑞与百时美施贵宝、信达与辉瑞的两笔大交易驱动。交易本身没有停滞，反而在加速。

如果只看股价，港股医药板块年初至今的表现确实疲弱。但如果看交易数据，中国药企的对外授权正在经历一个历史性的爆发期。

2025年全年，中国药企对外授权交易的前期付款总额达到了72.29亿美元，交易数量193笔，均创历史新高。2026

[... middle omitted ...]

和流动性因素，而不是产业逻辑的破坏。

4. 报告尚未完全回答的关键问题：竞争优势的“可迁移性”到底有多强？

尽管报告对中国药企的研发节点地位给出了强有力的论证，但它对一个更深层次的问题着墨不多：如果美国真的通过行政手段或税收激励，推动本土生物技术产业重建类似中国的早期研发基础设施，需要多长时间？

报告提到，中国过去五年建立了高效、经验丰富、技术积累深厚的研发生态系统。但这种生态的形成，既依赖于大量受过良好教育的研发人才，也依赖于临床资源的可及性、监管效率以及资本市场的支持。这些要素在美国是否可以被快速复制？如果可以，需要多少资本投入和时间成本？

从历史经验看，产业生态的重建通常以十年为周期。美国的生物技术产业虽然资本雄厚，但在早期临床阶段的“资产密度”上，短期内很难赶上中国。这意味着，即使政策落地，美国药企在过渡期内依然需要依赖中国资产。但“过渡期”的具体长度，以及中国药企能否在这段时间内完成从“资产输出”到“自主商业化”的跃迁，是报告没有完全展开的关键变量。

\subsection{[R015] 市场真正低估的不是AI需求，而是数据护城河的再定价}
{\small\color{refgray}归类：AI与科技：从算力军备到数据护城河，上游设备与封测受益}

这份由某外资投行发布的周度研究报告，表面上覆盖了从AI到油气服务的广泛话题，但贯穿其中的一条暗线值得所有产业决策者停下来思考：市场对AI的定价正在从“谁在花钱”转向“谁的数据不可替代”。

报告的核心判断并非来自某个单独的结论，而是来自一个被大多数投资者忽略的结构性信号——当欧洲媒体与互联网分析师Adam Berlin用一套新的AI框架重新评估公司时，得分最高的并非芯片公司，而是那些拥有专有、难以复制的数据库的企业。RELX在10分制中拿到9分，Pearson 8分，Springer Nature 7.5分。而这些公司恰恰被市场错误地归类为“AI风险敞口”标的。

这个信号意味着什么？AI的竞争正在从算力军备竞赛转向数据资产的深度变现。而这一转变，将对当前市场对AI产业链的定价逻辑产生根本性冲击。

1. 半导体集中的经济价值不可持续，数据层才是真正的价值洼地

报告中最值得反复咀嚼的判断来自Jim Covello：半导体领域的经济价值集中是不可持续的。这并非唱空芯片，而是在提醒一个被忽视的结构性问题——当前AI投资回报的最大瓶颈不在算力，而在数据结构和编排。

报告明确指出，尽管企业级AI投入巨大，但效果参差不齐。问题的核心在于，大多数企业还没有建立起能够真正释放AI价值的数据基础设施。这解释了为什么Adam Berlin的AI框架将“专有、难以复制的数据库”作为核心评分维度。

从投资角度看，这意味着当前市场对AI产业链的估值分布存在系统性偏差。半导体公司享受了AI叙事带来的绝大部分估值溢价，而真正掌握数据护城河的企业却被市场错误地贴上了“AI风险”的标签。这种定价错位的修复，可能成为未来12-18个月最重要的alpha来源之一。

2. 欧洲媒体与互联网的AI护城河比想象中更深——但市场还没有看懂

Adam Berlin对欧洲媒体与互联网公司的全面覆盖，提供了一个难得的观察窗口。他的核心论点是：AI对媒体和互联网公司的颠覆风险被高估了，因为这些公司的数据护城河比表面看起来更坚固。

具体来说，通用大语言模型无法实时抓取分类信息（如房产、汽车广告），DIY中介也始终是利基市场。这意味着Rightmove、AutoTrader和Scout24等平台的核心价值不会被AI轻易侵蚀。RELX在法律领域的数据可靠性更是难以复制的护城河——在需要高度准确性的场景下，用户不会接受“AI可能出错”的风险。

更值得关注的是，Berlin的买入评级与GS“AI风险篮子”高度重叠。这组矛盾数据本身就是一个重要的投资信号：市场对AI风险的定价可能出现了系统性的误判。当一个分析师用严谨的框架证明“AI风险”其实是“AI护城河”时，那些被错误归类的公司就存在明显的重估空间。

在AI叙事占据几乎所有注意力的当下，报告对油气服务公司TGS的分析提醒我们：市场焦点过于狭窄时，往往意味着某些结构性机会被系统性低估。

报告的核心逻辑有三层：第一，国际石油公司正在重新增加勘探支出，这一趋势在中东地缘事件之前就已经开始，其根本驱动力是储量寿命下降和页岩油成熟。第二，地震勘探行业已经完成整合，TGS等公司在等待上行周期的过程中获得了真正的定价权。第三，地震勘探支出的特点是“历史性的大额且集中”，一旦启动，股价重估速度极快。

这里的关键洞察是：这不是一个单纯的地缘政治交易，而是一个基于行业基本面的结构性机会。市场当前对能源服务的关注度极低，这种“无人问津”的状态本身就是一个有价值的信号。报告没有完全回答的问题是：这一轮上行周期的幅度和持续时间能否超过市场预期？这需要读者在完整报告中寻找更详细的供需平衡分析。

4. 西门子能源的结构性故事正在被低估——电网业务才是真正的引擎

报告新增的Conviction List标的中，西门子能源的案例提

[... middle omitted ...]

亿欧元。这个数字如果兑现，意味着公司不仅是一个增长故事，还是一个显著的价值回报故事。但这里的假设仍需验证——资本回报的规模和节奏取决于公司未来的现金流生成能力和管理层意愿。

5. 报告尚未完全回答的关键问题：AI投资的回报何时才能在企业层面显现？

尽管报告提供了丰富的分析框架，但有一个核心问题仍未得到充分回答：企业级AI投资何时才能从“成本中心”转化为“利润中心”？

报告提到企业AI效果“最多是混合的”，但没有深入分析为什么。这里有几个可能的解释方向：第一，企业数据基础设施的改造需要时间，而这一过程往往被低估；第二，AI应用的价值体现在流程优化而非收入增长，这使得ROI难以量化；第三，大多数企业还没有找到将AI能力转化为定价权或市场份额的路径。

这个问题对投资者的含义是：当前AI产业链的估值可能过度反映了“技术可行性”，而低估了“商业可行性”的难度。当市场开始关注后者时，那些无法证明AI投资回报的企业可能面临估值压缩，而那些拥有清晰变现路径的公司则会获得溢价。

\subsection{[R016] 中国科技股的核心机会不在下游，而在上游的设备与封测}
{\small\color{refgray}归类：AI与科技：从算力军备到数据护城河，上游设备与封测受益}

五月中国科技板块反弹18\%之后，市场对近期回调的看法分化。多数讨论集中在AI应用、消费电子复苏节奏和芯片设计公司的估值上。但某外资投行最新研报给出的判断与此不同：当前真正值得加仓的，不是下游的芯片设计或终端组装，而是上游的半导体设备、封测和代工。这个排序的逆转，背后是供应链利润分配的结构性变化。

这份研报将半导体子行业的优先顺序重新排列为：半导体设备大于封测，封测大于代工，代工大于无晶圆设计。对于长期关注中国科技股的投资者而言，这个排序的变动意味着什么？它是否暗示了市场共识的盲点？

简单说，上游企业的订单可见度正在从“季度”拉长到“两年”。下游企业的利润则面临成本上升和需求分化的双重挤压。这不是一个短期的交易信号，而是一个可能持续到2027年的结构性配置逻辑。

研报给出的核心数据是：2026年，本土WFE（晶圆前端设备）供应商来自存储和先进逻辑的需求，预计将分别同比增长超过50\%和30\%，且2027年还有进一步上行空间。这个判断的底层支撑，是长鑫存储和长江存储的IPO时间表推进——这两家企业的资本开支计划正在从“不确定”变为“可预期”。

更关键的是，这个增长周期与以往不同。过去中国半导体设备公司的订单多依赖单一客户的扩产节奏，波动大、可预测性差。但本轮驱动因素更加多元：存储扩产之外，本土AI芯片从2026年下半年开始进入量产爬坡期，将带动后端设备（先进封装、测试等）的增量需求。这意味着设备公司的收入来源正在从“单引擎”变为“双引擎”，增长的可预测性显著提升。

对于估值而言，这带来的不是简单的盈利上调，而是估值框架的切换——从“周期股”切换到“成长股”。研报明确提到，“延长增长可见性有助于支撑该板块的高估值”。这句话的实际含义是：市场此前对设备股的高估值存在疑虑，但如果订单能见度从12个月拉长到24个月，那么当前估值就有了基本面支撑。

2. 封测与代工正在享受一轮“量价齐升”，但利润分配并不均匀

研报观察到，2026年第二季度，代工厂和封测厂出现了“强劲且持续的涨价”。背后的逻辑是产能利用率提升，以及上游材料成本上涨的传导。这听起来像是整个板块的利好，但细看之下，利润分配并不均匀。

涨价能否转化为利润，取决于两个变量：折旧压力和产品结构。代工厂和封测厂正面临因资本开支增加而上升的折旧负担，这会侵蚀毛利。真正能跑赢的，是那些对AI相关产品（如GPU、PMIC）敞口更高的封测厂——这些产品的毛利率更高，涨价传导也更顺畅。

这意味着，投资者不能简单地把“封测涨价”等同于“所有封测股受益”。真正的区分度在于：哪些公司拿到了AI相关的封装订单，哪些公司还在做传统的消费电子封装。前者可以享受“量价齐升”的正向循环，后者可能面临“涨价但被折旧吃掉”的尴尬。

研报对下游的判断采取了“精选”策略，而非全面看多。从需求端看，云服务提供商的资本开支依然强劲，AI应用的渗透正在拉动服务器和AI芯片的采购。汽车和工业领域的组件需求也在2026年第二季度出现复苏。但消费电子——尤其是安卓手机和AIoT设备——的需求正在走弱。

这个分化意味着什么？对于本土GPU和AI ASIC公司而言，本土代工产能的增加正在缓解供应瓶颈，新一代芯片的量产有望加速销售增长。但对于那些高度依赖消费电子、且无法将上游涨价转嫁给客户的无晶圆设计公司而言，利润压力正在累积。

研报没有明说但逻辑上可以推导的是：下游的“精选”策略，本质上是在押注那些有定价权或产品不可替代性的公司。如果一家芯片设计公司的主要客户是手机厂商，且产品可替代性高，那么上游涨价和下游需求走弱的双重压力，可能会在2026年下半年集中体现。

研报对上游持乐观态度，但有一个关键变量没有被充分讨论：涨价的可持续性。当前封测和代工的涨价，部分原因是材料成本上涨的被动传导。如果材料成

[... middle omitted ...]

？这决定了估值能否从“周期”切换到“成长”。

第二层是涨价传导能力。封测和代工公司的涨价，是成本驱动的被动传导，还是产能紧张驱动的主动提价？如果是后者，利润改善的持续性更强。

第三层是产品结构。公司对AI相关产品的敞口有多大？AI封装、GPU测试等业务的毛利率是否显著高于传统业务？这决定了在行业整体涨价中，哪些公司能拿到更大的利润份额。

这个框架可以帮助投资者避免“买整个板块”的粗放策略，转而聚焦那些真正受益于结构性变化的公司。

第一个风险是IPO时间表的变数。研报将设备需求增长的重要假设建立在长鑫存储和长江存储的IPO推进上。但这两家公司的上市进程受到政策、市场环境和地缘政治等多重因素影响，时间表可能存在不确定性。如果IPO延迟，对应的资本开支节奏也可能低于预期。

第二个风险是AI芯片量产的实际爬坡速度。研报假设2026年下半年本土AI芯片开始量产，但芯片从流片到量产再到良率爬坡，中间存在大量技术风险。如果量产进度慢于预期，后端设备的需求增量也会相应推迟。

\subsection{[R017] GLP-1口服药的战场，IQVIA数据正在低估真实竞争烈度}
{\small\color{refgray}归类：医药与GLP-1：口服药渗透率加速，多肽制造绿色化成瓶颈}

过去两个月，GLP-1赛道最受关注的话题已经从“谁在减重效果上更优”转向“口服药能否真正打开增量市场”。每周处方数据追踪，看似枯燥的数字波动，背后却隐藏着两个关键判断：第一，市场正在用注射剂的逻辑来评估口服药，这可能导致对需求潜力的系统性低估；第二，由于IQVIA数据捕获率存在显著偏差，投资者对两款核心口服药——诺和诺德的Wegovy Pill和礼来的Foundayo——的竞争态势判断，可能正在偏离真实情况。

某外资投行最新发布的周度追踪报告，虽然标题只提到“长周末影响周处方量”，但仔细拆解后会发现，这份报告真正有价值的部分，不在于那6\%的周环比下降，而在于它揭示了一个尚未被充分定价的结构性变量：口服GLP-1的处方数据捕获率差异，正在扭曲市场对两款药物相对表现的认知。本文基于该报告的核心数据与逻辑框架，提炼出几个值得产业决策者和投资者认真对待的洞察。

1. 周度处方波动是噪音，真正重要的是口服药渗透率正在加速提升

截至5月29日当周，GLP-1类药物的总处方量为244万张，环比下降约6\%。报告明确指出，这一下降与长周末有关，属于预期内的季节性扰动。如果只看这一周的数据，很容易得出“市场增长趋缓”的错误结论。

真正值得关注的信号是：4周同比增速仍维持在39.2\%的高位，与上周的39.3\%几乎持平。这意味着，即便剔除短期波动，GLP-1的整体需求仍在以接近40\%的年化速度扩张。更关键的是，这一增长的主要驱动力正在从注射剂转向口服药。

从市场份额结构看，礼来整体（Zepbound、Mounjaro、Foundayo合计）在总处方量中的占比已达59.7\%，诺和诺德为40.3\%。但口服药的渗透率变化才是决定未来格局的关键变量。Wegovy Pill上市后第8周，周处方量已达13.4万张，虽然本周环比下降约8\%，但考虑到其仍处于早期放量阶段，这一绝对数字本身已经说明口服GLP-1的市场接受度远超此前预期。

对于产业决策者而言，这意味着：如果企业还在以“口服药只是注射剂的补充”来制定战略，很可能已经落后于市场实际节奏。口服GLP-1正在从“替代方案”变为“主流选择”。

2. IQVIA数据对Wegovy Pill的捕获率可能只有50\%，市场对诺和诺德的竞争力存在系统性低估

这是整份报告中最具冲击力的判断，也是大多数投资者容易忽略的技术细节。

报告引用了诺和诺德在2026年年度股东大会上的说法：Wegovy Pill在第10周的周处方量是替尔泊肽（假设为Zepbound）的3倍。但同期IQVIA数据显示，这一倍数仅为1.6倍。差距从何而来？核心在于IQVIA对Wegovy Pill的处方数据捕获率可能只有约50\%。

报告分析师对此给出了三重解释：诺和诺德拥有更及时、更细颗粒度的数据源；通过NovoCare等渠道开出的处方可能尚未被IQVIA完全收录；未来依从率的变化也会影响实际患者人数与处方量之间的对应关系。

更值得关注的是，报告对礼来口服药Foundayo的IQVIA捕获率估算约为70\%。这意味着，两款口服药在IQVIA系统中的数据偏差方向是一致的——都被低估，但Wegovy Pill被低估的程度更大。如果这一判断成立，那么市场基于IQVIA数据得出的“Wegovy Pill领先优势正在收窄”的结论，可能是错误的。

对于投资者而言，这一差异意味着：诺和诺德在口服GLP-1领域的真实竞争优势，可能比公开数据所呈现的要大得多。如果IQVIA捕获率问题在未来几个季度得到修正，诺和诺德的相对估值可能面临重新定价。

3. Foundayo的“慢启动”不是失败信号，而是全新分子上市的正常节奏

与Wegovy Pill的快速放量相比，礼来的Foundayo在第8周仅录得16,982张周处方量，环比增长6.1\%

[... middle omitted ...]

nd是注射剂，且当时市场对GLP-1的认知已较为成熟，这一对比并不公平。更合理的参照是：Foundayo作为一个全新口服分子，其早期处方曲线与Zepbound早期的斜率其实相当接近。

对于礼来的投资者而言，这里需要保持耐心。报告分析师认为，Foundayo的放量“需要数月而非数周”，这一判断在逻辑上是自洽的。真正需要关注的时间节点是第三季度DTC广告启动后的数据变化。

4. 新患者流入仍在加速，但Wegovy Pill对Zepbound NBRx的虹吸效应已开始显现

报告中最值得警惕的信号来自新品牌处方（NBRx）数据。虽然GLP-1类药物的总新患者启动量仍处于高位（截至5月22日当周约22.9万），但结构正在发生变化。

Exhibit 18显示，Zepbound的NBRx在Wegovy Pill上市后出现明显下滑。虽然报告强调，Zepbound的绝对NBRx数量今年仍维持在每周约6万张的水平，但从市场份额角度看，礼来的NBRx占比已从此前的约53\%有所回落。

\subsection{[R018] 保险股反弹背后，市场真正低估的是渠道结构重塑的定价权}
{\small\color{refgray}归类：保险与支付：渠道重塑与信任定价}

如果你在过去三个月看到中国保险公司的股价反弹，自然会问一个问题：这是跟着A股大盘走的贝塔，还是行业基本面发生了结构性改善？

某外资投行在5月底对中国主要保险公司进行了一轮深度调研，覆盖平安、太保、太平、众安、国寿、新华和人保财险。调研的结论可以浓缩为一句话：行业基本面确实在变好，但变好的方式比多数人想象的更微妙——不是单纯靠股市上涨拉高投资收益，而是渠道费用监管收紧正在悄悄改变竞争格局，让大型保险公司重新拿回定价权。

1. 银保渠道的“监管红利”正在重塑寿险竞争格局，大型险企是最大受益者

2025年3月，国家金融监管总局发布了银保渠道费用新规。这次调研中，所有寿险公司都对这一政策表达了欢迎态度。这不是客套话，而是有实质含义的。

过去几年，银保渠道的竞争本质上是费用战。中小保险公司为了抢占银行网点，不惜以高额手续费换取规模，导致整个渠道的利润率被压得很低。新规的核心在于对渠道费用进行更精细化的审查，这直接抬高了中小公司的竞争门槛。

大型险企的逻辑很清楚：银行在费用受控的条件下，会更倾向于与偿付能力充足、品牌认知度高、能够提供长期服务的大型保险公司合作。这不是猜测，而是已经在发生的事实。以中国太平为例，其银保渠道的VONB贡献在2025年约为35\%，但管理层明确表示这个比例有潜力逐步提升到50\%。太平人寿已经是银保渠道前七大寿险公司之一，这恰恰是监管整顿后的结果。

当然，监管红利不是免费的午餐。多家公司都承认，短期内销售激励的减少会造成一定的业务扰动。国寿在经历了2026年一季度VONB同比增长75\%的异常高增长后，管理层预计后续增速会回归正常，部分原因就是银保新规带来的短期冲击。但关键在于，这种冲击是一次性的，而竞争格局的改善是结构性的。

这意味着什么？对于投资者而言，银保渠道的利润率能否持续改善，取决于大型险企能否将节省下来的费用真正转化为利润率，而不是重新投入到价格战中去。这份调研中，保险公司普遍预期“大部分节省将让利给投保人”，这暗示利润率改善的幅度可能有限，但规模的集中化趋势是确定的。

2. 权益投资策略正在发生静默转变：从追逐弹性到配置“类债型”高股息

调研中一个容易被忽视但极其重要的信号，是保险公司对权益投资策略的一致表述。

几乎所有受访的寿险公司都强调了一个共同点：不会因为近期股市反弹而大幅提升权益配置比例。平安、国寿、新华都明确表示，整体权益配置比例将保持相对稳定。新华甚至指出，其2025年权益配置比例为22\%，一季度略有下降，且中期内不打算提高到25\%-30\%的区间。

这听起来像是一个保守的信号，但深入看会发现，真正重要的不是配置比例，而是配置结构的变化。

报告反复提到一个关键指标：OCI权益投资的占比在提升。OCI（其他综合收益）账户下的股权投资，其市值波动不计入当期损益，这使得保险公司可以将高股息股票作为“类债型”资产来持有，提供稳定的分红收益和久期匹配。国寿的OCI权益占比已达到约30\%，与行业平均水平大致相当。

这些保险公司对高股息股票的“目标股息率”普遍设定在4\%或以上。平安特别提到，如果负债成本继续下降，这个门槛还可以降低。这意味着，在利率长期下行的背景下，保险公司正在主动将权益资产“债券化”——不是追求股价弹性，而是追求稳定的现金流回报。

这个趋势对市场的含义是深远的。当中国最大的机构投资者群体开始系统性地将高股息股票作为长期配置核心，而不是交易性资产，这些股票的估值逻辑就会从“周期股”转向“类固收资产”。这或许可以解释为什么银行股和部分公用事业股在利率下行的环境中反而获得了重估。

但报告也留下了一个未解的问题：如果股市出现大幅回调，这些OCI账户中的高股息股票是否会因为浮亏而影响偿付能力？保险公司对此的应对机制，报告没有完全展开。

3. 代理人队伍

[... middle omitted ...]

品和健康险，仍然需要代理人的面对面服务。

这意味着，未来寿险行业的竞争将不再是“谁的人多”，而是“谁能真正渗透到中产及以上客群的保障需求”。太保在调研中提到，其代理人渠道的VONB利润率已接近30\%，主要由长期缴储蓄产品驱动。但银保渠道的利润率只有15\%左右，且主要通过趸缴产品实现。如果太保能够通过减少趸缴产品的占比来提升银保利润率，这将是一个重要的边际改善。

但这里存在一个尚未被充分讨论的风险：如果代理人数量真的稳定下来，而银保渠道的利润率提升有限，那么整个行业的VONB增长将越来越依赖规模扩张，而不是利润率提升。这恰恰是报告没有完全回答的问题——在费用监管红利释放完毕后，行业下一轮增长的动力来自哪里？

4. 财险的承保利润率改善有结构性支撑，但新能源车险仍是最大变量

财险板块的调研信号相对清晰。人保财险将自身优于同业的综合成本率（COR）归因于四个因素：更好的风险定价、内部费用管理、与主机厂的紧密关系、以及理赔管理的专业能力。这些优势不是短期内能够复制的。

\subsection{[R019] AI基础设施的供给瓶颈，正在从芯片设计延伸至整个物理供应链}
{\small\color{refgray}归类：AI与科技：从算力军备到数据护城河，上游设备与封测受益}

这份来自某外资投行的研报，在Broadcom 2Q26财报后第一时间更新了对AI半导体产业链的判断。表面上看，它是对一家公司业绩的追踪。但当我们把数字背后的信号串联起来，一个更重要的结构性叙事浮现出来：AI基础设施的竞争，已经从“谁的芯片算力更强”进入了“谁能锁定物理世界的供应”阶段。

报告中最值得关注的判断，不是Broadcom AI收入突破100亿美元的时间表，而是管理层披露的AI半导体订单积压已超过300亿美元，且订单能见度已延伸至2028年。客户之所以提前下单，不仅是为了锁定晶圆和HBM产能，更是为了确保电力基础设施的配套到位。这是一个信号：AI基础设施的供给瓶颈，正在从硅片层面蔓延到整个物理世界。

以下，我们从四个层次来拆解这份报告的核心洞察，并在最后提出一个报告尚未完全回答的关键问题。

1. 订单能见度延伸至2028年，供应链约束正在从晶圆向物理基础设施转移

报告中最硬核的数字，是Broadcom AI半导体季度订单积压超过300亿美元。管理层明确表示，客户正在提前下单，原因不仅仅是晶圆和HBM的交付周期长，更关键的是电力基础设施的配套周期。

这指向一个被市场低估的变化：AI基础设施的供给约束，正在从“芯片产能”向“数据中心电力容量”迁移。当电力成为瓶颈，客户不得不提前数年锁定芯片供应，以确保在拿到电力配额时，配套的计算设备已经就位。反过来，这也意味着芯片供应商的订单能见度被大幅拉长。

报告提到，这份研报的亚洲供应链调研确认，3至5年的长期协议已经越来越多地覆盖基板、HBM和CCL等关键瓶颈环节，而先进制程晶圆仍维持年度谈判周期。这隐含了一个判断：供应链约束的关键变量，正在从台积电的CoWoS产能向更上游的物理材料和电力转移。基板、HBM、CCL的产能扩张周期长于晶圆，一旦这些环节出现瓶颈，恢复时间将更长。

对投资者的含义：关注AI基础设施的“物理层”供应商——基板、HBM、CCL、封装测试——可能比关注芯片设计公司本身，更能捕捉到供给约束下的定价权溢价。报告明确列出的受益标的包括TSMC、ASE、SK Hynix、Samsung、Unimicron、Ibiden和EMC。

2. 客户组合多元化正在加速，但真正的赢家是那些能适应“半定制化”模式的供应商

Broadcom的客户名单正在快速扩容。除了Google、Meta、Anthropic、OpenAI这些已公开的客户，报告还基于供应链调研指出，另外两家客户——字节跳动和软银/ARM——将在2026年底开始出货，2027年加速放量。目前这些客户的采购订单金额已达约60亿美元。

但报告最值得细读的部分，是管理层首次公开承认Google将保持TPU供应链的多元化。这意味着，即使Broadcom与Google签有长期协议，Google仍将与联发科在TPU上合作，甚至可能与Marvell在LPU类芯片上合作。

报告指出，长期来看，大型云服务商同时与2至3家后端设计服务商合作，可能成为常态。驱动因素有两个：一是产品类型更加多样化（AI加速器、AI CPU、LPU类芯片等）；二是随着云服务商前端设计能力的成熟，它们会尝试不同的商业结构（半定制化vs.全定制化）。

这里的关键洞察是：联发科和Alchip这类毛利率较低（联发科40\%以上、Alchip十几\%）的供应商，反而可能成为赢家，因为云服务商在转向半定制化模式时，会更倾向于与利润率结构更“薄”的供应商合作。相比之下，那些同时提供前端设计服务的交钥匙供应商（毛利率60\%以上），在价格谈判中可能处于劣势。

对投资者的含义：不要简单地把“客户多元化”等同于“Broadcom被分走份额”。真正的竞争格局变化，是商业模式的分化——半定制化供应商的崛起，可能重新定义整个AI芯片设计服务市场的利润分

[... middle omitted ...]

太网交换机（Tomahawk 6已出货超过一年），再到1.6T DSP和CW/EML激光器等CPO构建模块，正在试图定义行业标准。下一代的200Tb交换机计划在本季度流片。

报告特别指出，随着agentic AI推动更大规模的扩展域和规模化域，高带宽数据传输的需求正在快速增长。但Broadcom对CPO（共封装光学）的推进节奏仍相对保守，预计其CPO量产时间表比NVIDIA晚约两年。

这里有一个隐含判断：网络互联的竞争，正在从“提供更高带宽”转向“定义互联标准”。谁能在交换机、SerDes、光学模块上确立标准，谁就能在AI基础设施的“互联税”中获得最大的份额。

对投资者的含义：市场目前对AI基础设施的关注仍集中在算力芯片上，但网络互联的竞争格局和定价权变化，可能在未来12至18个月内成为新的投资主线。CPO的渗透节奏、交换机带宽的升级周期、铜缆与光模块的替代关系，都是需要持续跟踪的变量。

4. 非AI半导体正在复苏，但驱动力并非传统需求，而是AI的“挤出效应”

\subsection{[R020] ASCO 2026第三天：市场低估的不是PD-1/VEGF组合的竞争，而是IL-2赛道“可成药性基准”的转移}
{\small\color{refgray}归类：医药与GLP-1：口服药渗透率加速，多肽制造绿色化成瓶颈}

ASCO 2026第三天：市场低估的不是PD-1/VEGF组合的竞争，而是IL-2赛道“可成药性基准”的转移

每年ASCO的第三天，往往是泡沫退潮、信号浮现的时刻。前两天的喧嚣过后，临床数据开始接受更冷静的拆解。

某外资投行在ASCO 2026第三天发布的研报，表面上是一份常规的会议纪要，但细读之下，它揭示了一个正在发生的结构性变化：中国创新药行业正在从“跟随式创新”进入“基准定义”阶段。这个判断，比任何单个分子的数据都更重要。

市场当前关注的焦点是PD-1/VEGF双抗的竞争格局，但报告真正值得重视的信号，是IL-2赛道正在经历的“可成药性基准”的转移。谁定义了基准，谁就掌握了下一轮定价权。

1. PD-1/VEGF的竞争已进入“数据细节决定估值分化”的阶段

HB0025在1L鳞状NSCLC中交出了ORR 84.5\%、mPFS 16.46个月的数据。单看数字，这组数据在同类分子中具有竞争力。但报告明确指出，存在几个需要谨慎解读的“caveats”：患者分期的偏倚、PD-L1表达水平的分布差异、以及低转移负荷患者的比例偏高。

这意味着什么？意味着PD-1/VEGF这个赛道的竞争，已经从“有没有疗效”进入了“疗效是否经得起亚组分析”的阶段。对于投资者而言，这不再是一个“买赛道”的逻辑，而是一个“挑分子”的逻辑。

AK112的Ph2数据作为参照系，HB0025的疗效信号需要放在同样的患者基线下去比较。报告没有给出结论，但暗示了方向：在鳞癌PD-L1 Personal reading notes and learning share only. Not investment advice.

PD-1/IL-2赛道，ANV600和AWT020都给出了验证数据，但βγ设计的IL-2分子依然受限于剂量窗口。IBI363作为可药性标杆，在2线免疫耐药EGFRwt非鳞NSCLC中，3mg剂量显示出OS优势，支持注册路径。

HB0025 20mg/kg Q3W联化疗，在1线NSCLC中展现高活性。鳞癌亚组尤其亮眼：ORR 84.5\%，mPFS 16.46个月，PD-L1<1\%和≥1\%患者的mPFS几乎一致（16.72 vs 16.46个月）。但需注意，与AK112的Ph2数据比，存在PD-L1、分期和低转移率的偏倚，G3+蛋白尿和TRAE死亡也是隐忧。

停药率仅1.2\%（化疗组11.2\%），中位剂量强度93.1\%。皮疹（17\%）和口腔炎（6.6\%）是主要减量原因。1线设计的核心争论点：单药vs联合vs SOC/维持，以及不同分子的临床表型差异。

China Biotechnology | Asia Pacific

PD-1/IL-2: ANV600/AWT020 add validation for PD-1-directed IL-2, but \$\textbackslash{}beta\textbackslash{}gamma\$ -design IL-2s still show dose window constraints; IBI363 remains the druggable benchmark. PD-(L)1/VEGF: HB0025 20mg/kg Q3W + chemo showed high 1L NSCLC activity, led by PD-L1-all-comer sqNSCLC efficacy signal. Caveats remain but bears watching. LBA04/HARMONi-6: refer to our takeaway note - China Biotechnology: ASCO 2026 Takeaways – HARMONi-6 Plenary. LBA05/RASolute302: 

[... middle omitted ...]

h sq/non-sq; ORR 84.5\%/65.6\%, mPFS 16.46/14.65mo, Sq PD-L1 Coverage Universe Investment Banking Clients (IBC) Other Material Investment ServicesClients (MISC) Stock RatingCategory Count \% of Total Count \% of Total IBC \% of RatingCategory Count \% of Total OtherMISC Overweight/Buy 1542 42\% 465 51\% 30\% 707 43\% Equal-weight/Hold 1571 43\% 369 40\% 23\% 723 44\% Not-Rated/Hold 3 0\% 0 0\% 0\% 1 0\% Underweight/Sell 551 15\% 86 9\% 16\% 201 12\% Total 3,

\subsection{[R021] 减肥药热潮的真正瓶颈不在需求，而在多肽制造的绿色化能力}
{\small\color{refgray}归类：医药与GLP-1：口服药渗透率加速，多肽制造绿色化成瓶颈}

市场对GLP-1类药物的关注几乎全部集中在需求端：全球肥胖药物支出预计在 年间以20\%以上的年复合增速增长，到2029年突破800亿美元，向千亿美元门槛逼近。但某外资投行最新发布的医疗健康周报揭示了一个被严重低估的结构性矛盾——多肽制造的可持续性瓶颈。这不是一个遥远的ESG议题，而是未来3-5年决定哪些CDMO能真正兑现产能扩张、哪些制药企业能锁定稳定供应链的核心变量。

报告的核心判断是：多肽API制造的高溶剂消耗和高过程质量强度（PMI），正在成为制约产能有效扩张的隐形天花板。 每生产1公斤GLP-1类多肽API，平均消耗约14吨溶剂，而典型小分子药物的这一数字仅为0.3吨。当行业从公斤级向多公吨级跃迁时，这种效率差距将直接转化为成本劣势、产能瓶颈和环境负债。

这不是一次技术路线的渐进改良，而是对多肽制造底层逻辑的重塑。那些率先掌握绿色合成工艺的CDMO，将在未来5年获得显著的竞争壁垒。

1. 多肽已成新药模态中规模最大、增速最快的赛道，但制造端的“脏活”被低估了

到2030年，多肽药物的市场规模预计将达到1200亿美元，年复合增长率18\%，在新型药物模态中规模最大、增速仅次于寡核苷酸和双抗。这一增长的绝对主力是GLP-1类药物——Semaglutide、Tirzepatide、Liraglutide三大品种已经将多肽从学术研究和化妆品原料的细分市场，推升为制药制造的核心赛道。

但报告明确指出了硬币的另一面：多肽API被美国化学会绿色化学研究所明确标记为“最不可持续的药物模态之一”。固体相多肽合成（SPPS）虽然是当前商业标准，但其过程质量指数（PMI）高达13,000以上——这意味着生产1公斤多肽API需要消耗13吨以上的原材料。对于年产量达到数公吨的GLP-1产品，这一数字意味着每年数万吨的溶剂消耗和废弃物处理压力。

这里的关键洞察在于：当市场只关注“产能扩张了多少升”时，真正决定有效产能的，是每升反应器能产出多少合格API。 高PMI不仅推高成本，更通过以下三种机制制约有效产能：第一，批次纯化中的大量中间产物被废弃，实际收率远低于理论值；第二，高溶剂消耗增加了废弃物处理成本和监管合规风险；第三，批次失败的概率随规模放大而上升，造成产能的不可预测性。

这意味着，单纯扩大反应器体积而不优化合成效率，本质上是在放大一个高浪费系统的规模。从投资角度看，衡量CDMO竞争力的核心指标应从“产能升数”转向“单位产能的PMI改善幅度”。

报告以Bachem为案例，详细拆解了多柱逆流溶剂梯度纯化（MCSGP）技术如何改变纯化环节的经济性。这是整份报告最值得深挖的技术细节之一。

传统批次纯化的核心问题是“中间产物浪费”：在分批操作中，色谱柱之间的过渡段含有大量纯度不达标的中间产物，这些产物通常被丢弃或需要反复处理。MCSGP通过两根连接的色谱柱连续运行，将这部分中间产物反复回传至柱内，直到达到99.3\%以上的纯度标准。

报告提供的量化数据极具说服力：相比传统批次纯化，MCSGP将溶剂消耗降低约30\%，总物料使用量降低70\%。更关键的是，原本需要5天的批次纯化流程被压缩到24小时连续运行。这意味着：第一，在制品控制检查次数大幅减少；第二，冻干体积显著缩小，释放QC资源；第三，设备利用率成倍提升。

从产业含义来看，这项技术的战略价值不仅在于降本。当CDMO能够在更短时间内完成更高纯度的纯化，实际上是在不增加反应器容量的前提下，创造了更多“有效产能”。对于面临产能瓶颈的GLP-1供应链，这种效率提升比新建产线更具性价比。报告没有完全展开的一个问题是：MCSGP的规模化部署是否存在技术天花板？从ETH Zurich的实验室概念到Bachem的商业化落地，这一技术目前仍处于早期推广阶段，其在不同序列长度和疏水性

[... middle omitted ...]

的ESG披露。

报告敏锐地指出了这一趋势对客户粘性的影响：当制药企业面临越来越严格的ESG合规要求时，能够提供量化PMI改善和溶剂回收数据的CDMO，将获得更稳定、更长期的合同。这不仅仅是成本竞争，而是从“价格-质量”二元竞争转向“价格-质量-可持续性”三维竞争。

但这里有一个值得追问的问题：PMI优化是否存在边际递减效应？当WuXi将PMI从13,000降至3,600后，进一步降低到1,000甚至更低的技术可行性和经济性如何？报告没有深入探讨这一极限，但这对评估CDMO的长期竞争优势至关重要。

4. 时间压缩能力正在成为CDMO的差异化武器，但“快”不等于“好”

Asymchem的Mazdutide案例展示了另一种竞争维度：时间。典型的长链多肽和GLP-1类药物的过程性能确认（PPQ）通常需要18-24个月，而Asymchem在12个月内完成了预PPQ和PPQ，6个月内提交NDA档案，6个月内通过动态现场审计，最终交付了超过5公斤、纯度大于98.5\%的批次。

\subsection{[R022] 市场真正低估的不是卡组织费率，而是“信任即服务”的定价权}
{\small\color{refgray}归类：保险与支付：渠道重塑与信任定价}

关于银行卡收单费率（Merchant Discount Rate, MDR）的讨论，几乎每隔一段时间就会在支付行业掀起一轮争论。反对者认为，2.4\%的信用卡费率是美国商户的沉重负担，稳定币、即时支付等替代方案将终结卡组织的“超额利润”。支持者则强调，卡支付带来的争议处理、欺诈保护、全球通用性等增值服务，远非简单的资金转移可比。

这份来自某外资投行的周末技术简报，用一个看似老生常谈的费率问题，揭示了支付行业最核心却最容易被误解的结构性事实：市场对卡组织费率的焦虑，本质上混淆了“资金转移”与“信任服务”两种完全不同的价值定价。而真正决定未来十年支付格局的，不是费率高低，而是谁能把“信任”这个抽象资产，转化为可规模化的商业模型。

1. 费率对比中最常被忽略的前提：借记卡与信用卡根本是两种商品

许多关于卡支付“昂贵”的批评，从一开始就站错了比较基准。该报告明确指出，借记卡与信用卡在功能上完全是两回事。借记卡对应的是“你已有的钱”，是资金转移；信用卡对应的是“你还没有的钱”，是信用延伸。

这一区分至关重要。当人们拿“银行直付”（pay-by-bank）或稳定币与信用卡比较时，实质是在拿一种资金转移工具，去对标一种集成了信用、奖励、争议处理、欺诈保障的复合服务产品。数据显示，美国受监管的借记卡平均MDR仅为0.7\%，远低于信用卡的2.4\%。即便加上不受监管的借记卡部分（约占全美借记卡交易量的三分之一），其费率也仅在1.7\%左右，仍然低于信用卡。

在欧洲，这一对比更加鲜明。欧盟对借记卡和信用卡的交换费（interchange fee）分别设有约20个基点和30个基点的上限，卡支付的总体费率在全球属于最低水平之一。但报告同样指出，欧洲的卡支付渗透率并未因费率受限而停滞——自2015年交换费上限实施后，欧洲的卡支付渗透率反而从33\%跃升至2025年的77\%。

这意味着什么？费率本身并不是决定支付工具普及的唯一变量。真正驱动用户和商户采用某种支付方式的，是它能否在“转移资金”这个最小公分母之上，构建出足够多的价值附加层。借记卡已经足够便宜，真正需要被审视的不是卡支付的整体成本，而是信用卡这个“信用+服务”复合品的定价是否合理。

2. 信用卡费率背后是数十年的信任基础设施，稳定币“便宜”的叙事经不起推敲

将信用卡MDR拆解开来，报告给出了一个令人清醒的构成分析。以美国信用卡2.8\%的MDR为例，其中网络转接与处理费约占1个百分点，欺诈成本1.5个百分点，争议处理1.2个百分点，而发卡奖励、获客与维护成本高达3个百分点。此外还有约2.5个百分点的资金成本。

简单说，商户支付的每一笔信用卡手续费中，超过一半被用于覆盖欺诈、争议和奖励——这些都是“信任”和“服务”的成本，而不是资金转移的成本。稳定币或银行直付之所以能在某些场景下宣称“更便宜”，恰恰是因为它们剥离了这些成本项。稳定币交易不可撤销，没有争议处理机制；奖励机制依赖的浮存金收益存在数学上限；而欺诈防范和KYC（了解你的客户）更是需要数年积累的治理能力。

报告中的一个数据点尤为尖锐：在Shopify平台上，标准线上卡支付费率为2.9\%加30美分，而USDC（美元稳定币）的支付费率同样是2.9\%加30美分。对于小型商户而言，稳定币并未带来任何成本优势。

这引出了一个更深层的结论：支付行业的真正护城河不在于技术，而在于治理。 资金转移本身是商品化服务——ACH（自动清算系统）的单笔交易成本不到1美分，比大多数稳定币还低。但信任和治理不是商品。卡组织用数十年时间，在全球超过1亿个商户网点和70亿张经过KYC的凭证之间，建立了一套可执行的规则体系。这套体系无法被单一的技术方案替代，因为它本质上是一套社会契约——商户接受高费率，换来的是更高的转化率、更低的欺诈损失、以及

[... middle omitted ...]

的调查数据，截至2025年1月，已有34\%的美国小型商户开始向信用卡用户收取附加费。更值得关注的是，正在审理中的MDL（Multi-District Litigation）和解协议如果获批，将进一步放宽商户收取附加费的限制。该协议还计划在未来五年内将信用卡交换费降低约10个基点，并对标准消费者信用卡费率设定上限。

这里需要特别指出的是，附加费仅适用于信用卡，不适用于借记卡——因为借记卡已经足够便宜。报告对此表达了谨慎的担忧：虽然目前附加费的渗透率仍然较低，但一旦MDL和解协议落地，附加费的普及速度可能加快。而商户收单机构（如Toast）本身也有动力推动附加费，因为它能带来ARPU（每用户平均收入）的提升。

然而，真正值得关注的不是附加费本身，而是它对用户体验的长期侵蚀。如果消费者在越来越多的场景中被要求为信用卡支付额外费用，卡支付的“无缝体验”优势将被削弱。这对卡组织而言是一个必须与发卡行共同应对的棘手问题：发卡行希望维持高交换费，但商户和消费者的接受度存在边界。

\subsection{[R023] 日本汽车业IT投资正在踩刹车，但这恰恰暴露了另一条增长主线}
{\small\color{refgray}归类：欧洲汽车与日本IT：结构性变化下的赢家与输家}

一份来自某外资投行的研报，通过对三家非覆盖公司（Systena、Argo Graphics、DIT）的实地调研，揭示了一个正在发生的结构性变化：日本汽车行业的IT投资意愿正在快速降温。但比这更重要的是，这份报告同时指向了一个被市场低估的替代增长引擎——金融行业的核心系统更新需求。

这不是一份关于“日本IT服务行业整体好不好”的报告。它的真正价值在于，它用三家公司的微观数据，拆解了“谁在收缩、谁在扩张、以及收缩与扩张之间的传导逻辑”。对于跟踪日本IT服务板块的投资者来说，这份报告提供了一个关键的观察框架：在汽车行业这个传统大客户的IT支出放缓时，哪些公司有能力完成切换，哪些公司会被困在原地。

1. 汽车行业的IT投资收缩不是偶发现象，而是三个信号叠加的结果

研报中三家公司给出的信息高度一致，但程度不同。Systena说“部分整车厂的投资意愿正在减弱”；Argo Graphics直接给出了六年来首次营业利润同比下降的指引，原因是“主要汽车客户的投资受到抑制”；DIT的情况更严峻，其嵌入式解决方案业务的销售额和毛利润在3Q均出现同比下降，且公司明确表示“主要汽车客户的投资意愿正在迅速下降”。

这不是一个公司的个案，而是一个行业层面的信号。三个公司分别覆盖了车载系统集成（Systena、DIT）和汽车CAD/半导体相关系统（Argo Graphics），它们的共同指向是：日本汽车行业正在经历一轮IT投资的周期性回调。

这意味着什么？对于投资者而言，需要重新审视那些汽车行业收入占比较高的IT服务公司的盈利预期。虽然报告指出目前没有大型系统集成商对汽车行业的销售敞口特别高，但这恰恰是风险所在——当风险分散在多个子行业时，市场往往低估了它集中爆发的可能性。如果汽车行业的投资收缩持续到FY6/27，那么那些看似“多元化”的公司，其汽车业务的下滑可能需要其他板块的超额增长来弥补，而这并非总能实现。

2. 金融行业的核心系统更新正在成为新的需求支柱，且趋势具备持续性

与汽车行业的收缩形成鲜明对比的，是Systena的数字集成业务表现。该业务的订单持续强劲，特别是金融领域的核心系统更新项目，咨询量处于高位。Systena管理层明确表示，基于当前的订单趋势，这一增长势头有望延续。

这一点值得深入思考。金融行业的核心系统更新与汽车行业的SDV（软件定义汽车）投资，在本质上有一个关键区别：前者是监管驱动和系统老化驱动的刚性需求，后者是战略投资驱动的弹性需求。刚性需求的韧性更强，受宏观经济波动的影响更小。

报告没有展开讨论的是，日本金融机构的系统更新需求背后，有几个结构性因素在支撑：一是日本银行业长期以来的老旧系统（legacy system）问题，二是监管对系统安全性和数据管理的要求持续提升，三是数字化转型带来的竞争压力。这些因素叠加，使得金融IT支出的增长更具确定性。

对于投资者来说，这意味着在汽车行业IT投资放缓的环境下，那些在金融行业有深厚布局的系统集成商，可能会展现出更强的盈利韧性。Systena的数字集成业务营业利润率已经达到23.8\%，且仍在提升，这本身就是最好的证明。

3. 半导体领域的投资复苏正在分化，但区域集中度是一个被忽视的变量

Argo Graphics的案例提供了一个有趣的观察窗口。尽管汽车客户的IT投资在收缩，但半导体领域的投资正在出现复苏迹象。公司明确提到，多个主要客户计划在北海道、九州、三重等地的工厂扩大投资。

这里的关键洞察在于区域集中度。半导体工厂的建设周期长、投资规模大，而且一旦启动，相关的CAD、PLM（产品生命周期管理）和HPC（高性能计算）系统需求就会形成持续的订单流。对于像Argo Graphics这样在半导体领域有布局的公司来说，这提供了一个对冲汽车行业下滑的缓冲垫。

[... middle omitted ...]

原因是股权激励成本增加了8.4亿日元。Argo Graphics的利润下滑部分原因则是数据中心折旧成本上升。DIT则提到，为了达成利润指引，需要压缩战略投资（安全产品开发、AI投资、内部系统投资等）和奖金。

这些成本端的压力，报告只是点到为止，没有深入展开。但这是一个关键问题：在收入端面临汽车行业逆风的情况下，日本IT服务公司是否还有足够的空间来消化成本上升？如果成本压力持续，而收入增长放缓，那么利润率的压缩可能比市场预期的更严重。

更值得思考的是，DIT压缩战略投资的做法，虽然短期内保住了利润，但长期来看可能会削弱竞争力。在AI和安全产品上的投资不足，意味着它可能在下一轮技术周期中落后于竞争对手。这是一个典型的“短期利润 vs 长期竞争力”的取舍问题，而市场往往对这类取舍的定价不够充分。

对于读者来说，在跟踪这些公司时，不能只看收入和利润的绝对数字，更要关注成本结构的变化和投资决策的取舍。那些能够在成本压力下仍保持战略投入的公司，才更有可能在下一个周期中胜出。

\section{Disclaimer}
{\scriptsize\color{refgray}
This community is a paid learning and information-sharing community created for general educational purposes only. The membership fee is charged solely for access to community discussions, educational content, organization, and operational costs. It is not an advisory fee, management fee, performance fee, brokerage fee, or compensation for personalized financial, investment, legal, tax, or professional advice.\par
I participate and share information solely as an individual and for educational discussion among community members. I am not acting as a financial adviser, investment adviser, broker, dealer, portfolio manager, legal adviser, tax adviser, accountant, fiduciary, or any other licensed professional adviser. Nothing shared in this community constitutes, and should not be interpreted as, financial advice, investment advice, legal advice, tax advice, a recommendation, solicitation, offer, endorsement, instruction, or guarantee to buy, sell, hold, short, trade, or invest in any security, cryptocurrency, fund, derivative, strategy, or other financial product.\par
All content shared in this community, including messages, discussions, links, files, charts, examples, market commentary, watchlists, AI-generated or AI-polished summaries, and third-party materials, is general, impersonal, and educational in nature. It does not take into account any individual's financial situation, investment objectives, risk tolerance, investment horizon, liquidity needs, tax status, legal restrictions, or suitability requirements. No content should be relied upon as suitable for any specific person, account, portfolio, or financial decision.\par
Information shared in this community may be incomplete, inaccurate, outdated, speculative, AI-generated, AI-edited, or based on third-party internet sources that have not been independently verified. I make no representation or warranty as to the accuracy, completeness, timeliness, reliability, or suitability of any information shared.\par
Financial markets involve substantial risk, including the possible loss of principal. Past performance, historical data, hypothetical examples, simulated results, backtests, personal opinions, or educational examples do not guarantee future results. Each member is solely responsible for conducting their own research, exercising independent judgment, and making their own decisions. Before making any financial, investment, legal, or tax decision, members should consult qualified and licensed professionals.\par
Participation in this community, payment of any membership fee, or communication with me or other members does not create any adviser-client, fiduciary, professional, contractual, agency, partnership, employment, or other advisory relationship. By joining or participating in this community, you acknowledge and agree that you are solely responsible for your own decisions, actions, transactions, profits, losses, risks, and consequences, and that I shall not be liable for any direct, indirect, incidental, consequential, financial, legal, tax, or other loss or damage arising from or related to any content, discussion, information, or material shared in this community.\par
}
\end{document}
